\subsection{普遍可测函数演算}

在本节中，我们所考虑的希尔伯特空间都是复的。

设 u 是希尔伯特空间 E 上的正规部分算子。我们用 \(\mu_{x,y}\) （相应地，\(\mu_{y}\) ）表示 \((x,y)\in\mathrm{E}\times\mathrm{E}\) （相应地，y）关于 u 的谱测度。

设 \(y \in \mathrm{E}\) 。映射 \(\mathcal{K}(\mathrm{Sp}(u)) \to \mathrm{E}\) 由 \(f \mapsto f(u)(y)\) 所定义，它满足

\[
\| f (u) (y) \| ^ {2} = \int_ {\mathrm{Sp} (u)} | f | ^ {2} \mu_ {y} = \| f \| _ {\mathcal {L} ^ {2} (\mathrm{Sp} (u), \mu_ {y})} ^ {2}.
\]

因此，存在唯一的等距线性映射 \(ev_{y}\) ，它从空间 \(\mathrm{L}^{2}(\mathrm{Sp}(u),\mu_{y})\) 映入 E，使得 \(\mathrm{ev}_{y}(\widetilde{f})=f(u)(y)\) ，这里 \(\widetilde{f}\) 是函数 \(f\in\mathcal{K}(\mathrm{Sp}(u))\) 的类。

设 \(f\) 是定义在 \(u\) 的谱上的普遍可测函数。我们用 \(\mathrm{D}_{f}\) 表示这样的元素 \(y \in \mathrm{E}\) 所成的集：\(f\) 属于 \(\mathcal{L}^{2}(\mathrm{Sp}(u), \mu_{y})\) 。

命题 3. --- 设 f 是定义在 u 的谱上的普遍可测函数。集合 D \(_{f}\) 是 E 的稠密向量子空间。映射 y ↦ ev \(_{y}\) (f) 是 E 上的正规部分算子，其定义域为 D \(_{f}\) ，记作 f(u)。

对每个 \(x \in \mathrm{E}\) 及每个 \(y \in \mathrm{D}_f\) ，我们有 \(f \in \mathcal{L}^1(\mathrm{Sp}(u), \mu_{x,y})\) ，并且

\[
\langle x \mid f (u) y \rangle = \int_ {\operatorname{Sp} (u)} f \mu_ {x, y}.\tag{6}
\]

我们将证明下列更精确的陈述：

命题 4. --- 设 f 是 u 的谱上的普遍可测函数。

\begin{enumerate}
\def\labelenumi{\alph{enumi})}
\item
  集合 \(D_{f}\) 是 E 的稠密向量子空间。映射 \(f(u): y \mapsto \operatorname{ev}_{y}(f)\) 从 \(D_{f}\) 到 E 中是线性的，且当 f = 1 时与 \(1_{E}\) 重合；
\item
  设 \(g \in \mathcal{L}_{\mathrm{u}}(\mathrm{Sp}(u))\) 。对每个 \(y \in \mathrm{D}_f\) 及每个 \(x \in \mathrm{D}_g\) ，我们有 \(fg \in \mathcal{L}^1(\mathrm{Sp}(u), \mu_{x,y})\) ，并且
\end{enumerate}

\[
\langle g (u) x \mid f (u) y \rangle = \int_ {\mathrm{Sp} (u)} \overline {{g}} f \mu_ {x, y}.\tag{7}
\]

\begin{enumerate}
\def\labelenumi{\alph{enumi})}
\setcounter{enumi}{2}
\tightlist
\item
  假设 \(E = L^{2}(X, \mu)\) ，其中 X 是局部紧拓扑空间，\(\mu\) 是 X 上的正测度；又假设 \(u = m_{h}\) ，其中 \(h: X \to \mathbf C\) 是 \(\mu\)-可测的。我们有 \(f(m_{h}) = m_{f \circ h}\) 。
\end{enumerate}

由 IV, p. 266 的定理 1，我们可以假设处于断言 \(c\) ) 的情形，即 \(\mathrm{E} = \mathrm{L}^2 (\mathrm{X},\mu)\) 且 \(u = m_h\) ，其中 X 是局部紧拓扑空间，\(\mu\) 是 X 上的正测度，而 \(h\colon \mathrm{X}\to \mathbf{C}\) 是 \(\mu\)-可测的。我们将用 \(\widetilde{\varphi}\) 表示 \(\mathrm{L}^2 (\mathrm{X},\mu)\) 中函数 \(\varphi \in \mathcal{L}^2 (\mathrm{X},\mu)\) 的类。

设 \(\mathrm{S} = \mathrm{Sp}(m_h)\) 。我们用 \(\mu_{\widetilde{\varphi}_1, \widetilde{\varphi}_2}\) （相应地，\(\mu_{\widetilde{\varphi}}\) ）表示 \((\widetilde{\varphi}_1, \widetilde{\varphi}_2)\) （相应地，\(\widetilde{\varphi}\) ）关于 \(m_h\) 的谱测度，对所有 \((\varphi_1, \varphi_2) \in \mathcal{L}^2(\mathrm{X}, \mu)^2\) （相应地，每个 \(\varphi \in \mathcal{L}^2(\mathrm{X}, \mu)\) ）。

设 \(\varphi\in\mathcal{L}^{2}(X,\mu)\) 。谱测度 \(\mu_{\widetilde{\varphi}}\) 等于 S 上的像测度 \(h(|\varphi|^{2}\cdot\mu)\) （IV, p. 269 的引理 4）。由于函数 \(f\) 是 \(\mu_{\varphi}\)-可测的，我们有 \(\varphi\in D_{f}\) 当且仅当积分

\[
\int_ {\mathrm{S}} ^ {*} | f | ^ {2} d \mu_ {\widetilde {\varphi}} = \int_ {\mathrm{S}} ^ {*} | f | ^ {2} h (| \varphi | ^ {2} d \mu) = \int_ {\mathrm{X}} ^ {*} | f \circ h | ^ {2} | \varphi | ^ {2} d \mu
\]

都是有限的（INT, V, p. 70, § 6, 第 2 目, 命题 2）。这表明 \(D_{f}\) 是乘法算子 \(m_{f\circ h}\) 的定义域，而该定义域是 E 的稠密子空间（IV, p. 232 的命题 5）。

\(\mathrm{ev}_{\widetilde{\varphi}}\) 在函数类 \(g \in \mathcal{K}(\mathrm{S})\) 上的限制是把类 \(\widetilde{g}\) 对应到元素 \(g(m_h)(\widetilde{\varphi}) = m_{g \circ h}(\widetilde{\varphi})\) 的映射（IV, p. 269 的引理 4）。因此，映射 \(\mathrm{ev}_{\widetilde{\varphi}}\) 与从 \(\mathrm{L}^2(\mathrm{S}, \mu_{\widetilde{\varphi}})\) 到 \(\mathrm{L}^2(\mathrm{X}, \mu)\) 中的等距映射重合，后者是由映射 \(g \mapsto (g \circ h) \cdot \varphi\) （从 \(\mathcal{L}^2(\mathrm{S}, \mu_{\widetilde{\varphi}})\) 到 \(\mathcal{L}^2(\mathrm{X}, \mu)\) 中）取商得到的。

特别地，我们因此有 \(\mathrm{ev}_{\widetilde{\varphi}}(f)=m_{f\circ h}(\widetilde{\varphi})\) ，于是 \(f(m_{h})\) 作为从 \(D_{f}\) 到 E 中的映射与部分算子 \(m_{f\circ h}\) 重合。这就证明了断言 c)；若 f=1，则得到 \(\mathrm{ev}_{\widetilde{\varphi}}(1)=\widetilde{\varphi}\) ，这也就完成了 a) 的证明。

我们来证明断言 \(b\) )。设 \(\varphi_{1}\) 与 \(\varphi_{2}\) 是使 \(\widetilde{\varphi}_{1} \in \mathrm{D}_{f}\) 且 \(\widetilde{\varphi}_{2} \in \mathrm{D}_{g}\) 的函数。由（INT, V, loc. cit.）可得

\[
\begin{array}{l} \int_ {\mathrm{S}} ^ {*} | f g |   | \mu_ {\widetilde {\varphi_ {1}}, \widetilde {\varphi_ {2}}} | = \int_ {\mathrm{X}} ^ {*} | (f \circ h) (g \circ h) \varphi_ {1} \varphi_ {2} | d \mu \\ \qquad \leqslant \| (f \circ h) \varphi_ {1} \| \| (g \circ h) \varphi_ {2} \| \end{array}
\]

它是有限的，因此 \(fg \in \mathcal{L}^1(\mathrm{S}, \mu_{\widetilde{\varphi}_1, \widetilde{\varphi}_2})\) 。此外，这时我们有

\[
\begin{array}{c} \langle g (m _ {h}) (\widetilde {\varphi} _ {2})   |   f (m _ {h}) (\widetilde {\varphi} _ {1}) \rangle = \int_ {\mathrm{X}} \overline {{(g \circ h)   \varphi_ {2}}}    (f \circ h)   \varphi_ {1}    d \mu \\ = \int_ {\mathrm{S}} \overline {{g}}   f    d \mu_ {\widetilde {\varphi} _ {1}, \widetilde {\varphi} _ {2}}, \end{array}
\]

证明完毕。

定义 5. --- 映射 \(f \mapsto f(u)\) 由 IV, p. 271 的命题 3 所定义，它从 \(\mathcal{L}_{\mathrm{u}}(\mathrm{Sp}(u))\) 映到 E 上正规部分算子所成的集，称为与 \(u\) 相伴的普遍可测函数演算映射。

更一般地，设 T 是一个集合，并设 \(g\colon\operatorname{Sp}(u)\times\operatorname{T}\to\mathbf{C}\) 是一个映射；设 \(t\in T\) 使得函数 \(g_{t}\colon z\mapsto g(z,t)\) 在 \(\operatorname{Sp}(u)\) 上普遍可测。这时我们记 \(g(u,t)=g_{t}(u)\) 。

公式 (7) 特别地蕴含

\[
\| f (u) (y) \| ^ {2} = \int_ {\operatorname{Sp} (u)} | f | ^ {2} \mu_ {y}\tag{8}
\]

对所有 \(f \in \mathcal{L}_{\mathrm{u}}(\mathrm{Sp}(u))\) 及每个 \(y \in \mathrm{D}_f\) 成立。由此，取 \(f = 1\) ，我们推出 \(\mu_y\) 是总质量为 \(\| y\|^2\) 的正测度。

若 \(u\) 是希尔伯特空间 E 上的正规部分算子且 \(f \in \mathcal{K}(\mathrm{Sp}(u))\) ，则 \(\mathrm{D}_f = \mathrm{E}\) ，且部分算子 \(f(u)\) 这时是连续的，因为

\[
\| f (u) y \| \leqslant \| f \| _ {\infty} \| y \|
\]

对所有 \(y \in E\) 成立（依据 (8)）。自同态 \(f(u)\) 与 IV, p. 268 的定义 3 中的那个自同态重合。

推论 1. --- a) 对每个 \(f \in \mathcal{L}_{\mathrm{u}}(\mathrm{Sp}(u))\) ，部分算子 \(f(u)\) 是正规的，且 \(f(u)^{*} = \overline{f}(u)\) 。此外，\(f(u)\) 在 \(f \geqslant 0\) 时是正的，在 \(f\) 取实值时是自伴的；

\begin{enumerate}
\def\labelenumi{\alph{enumi})}
\setcounter{enumi}{1}
\item
  设 \(k \in \mathbf{N}\) ，令 \(f(z) = z^{k}\) ，\(z \in \operatorname{Sp}(u)\) 。我们有 \(f(u) = u^{k}\) ；
\item
  设 \(\lambda \in \mathbf{C} - \mathrm{Sp}(u)\) ，令 \(f(z) = (\lambda - z)^{-1}\) ，\(z \in \mathrm{Sp}(u)\) 。我们有 \(f(u) = \mathrm{R}(u, \lambda)\) ；
\item
  我们有 \(\beta(u) = b(u)\) ；
\item
  设 \(f \in \mathcal{L}_{\mathrm{u}}(\mathrm{Sp}(u))\) 与 \(g \in \mathcal{L}_{\mathrm{u}}(\mathrm{Sp}(u))\) 满足 \(|g| \leqslant 1 + |f|\) 。\(g(u)\) 的定义域是 \(f(u)\) 的核。
\end{enumerate}

鉴于谱定理（IV, p. 266 的定理 1），这由前面的命题分别结合以下各项而得：

\begin{enumerate}
\def\labelenumi{\alph{enumi})}
\item
  IV, p. 266 的引理 2, a)、IV, p. 253 的命题 23 及其推论；
\item
  IV, p. 255 的命题 24, b)；
\item
  IV, p. 252 的命题 22, b)；
\item
  IV, p. 266 的引理 2, a)；
\item
  IV, p. 232 的命题 6, b)。
\end{enumerate}

注. --- 对 \(f = \operatorname{Id}_{\operatorname{Sp}(u)}\) ，我们有 \(f(u) = u\) （断言 b) 取 \(k = 1\) ）。因此，u 的定义域与这样的 \(x \in E\) 所成的集重合：\(\operatorname{Sp}(u)\) 的恒等映射属于 \(\mathcal{L}^{2}(\operatorname{Sp}(u), \mu_{x})\) ；它特别包含这样的元素 \(x \in E\) ：测度 \(\mu_{x}\) 具有紧支撑。

下面的推论推广了 IV, p. 247 的命题 16 的推论以及 IV, p. 248 的命题 17。

推论 2. --- 设 u 是希尔伯特空间 E 上的正规部分算子。假设 E 非零。

\begin{enumerate}
\def\labelenumi{\alph{enumi})}
\item
  u 的谱非空；
\item
  对每个 \(\lambda \in \mathbf{C} - \mathrm{Sp}(u)\) ，预解式 \(\mathrm{R}(u, \lambda)\) 的范数等于 \(1/\delta\) ，其中 \(\delta > 0\) 是在 \(\mathbf{C}\) 中从 \(\lambda\) 到 \(u\) 的谱的距离。
\item
  对每个 \(\varepsilon > 0\) ，\(\varepsilon\)-伪谱 \(\mathrm{PSp}_{\varepsilon}(u)\) 是由 \(\lambda \in \mathbf{C}\) （距离 \(< \varepsilon\) ，从 \(\mathrm{Sp}(u)\) 量起）所成的集。
\end{enumerate}

假设 \(u\) 的谱为空。那么 \(u\) 是单射的，且自同态 \(u^{-1} = -\mathrm{R}(u,0)\) 是正规的（推论 1, c) 与 a)）。由于 \(\mathrm{Sp}(u^{-1}) \subset \{0\}\) （命题 15, a)），我们有 \(\mathrm{Sp}(u^{-1}) = \{0\}\) （I, p. 26, 推论 1）。由于 \(u^{-1}\) 是正规的，这就蕴含 \(u^{-1} = 0\) （I, p. 110, 例 1），而这是矛盾，因为 E 不为零。

为证明 \(b\) )，可以假设 \(u\) 是乘法算子 \(m_g\) ，作用于 \(\mathrm{L}^2 (\mathrm{X},\mu)\) ，其中 \(g\) 是赋以正测度 \(\mu\) 的局部紧拓扑空间 X 上的连续函数（IV, p. 266 的定理 1）。设 \(\lambda \in \mathbf{C} - \mathrm{Sp}(u)\) ，并设 \(\delta >0\) 是 \(\lambda\) 到 \(u\) 的谱的距离。为证明 \(\| \mathrm{R}(u,\lambda)\| = \delta^{-1}\) ，我们归结到 \(\lambda = 0\) 的情形，办法是把 \(u\) 换成 \(u - \lambda 1_{\mathrm{E}}\) 。于是实数 \(\delta\) 是 0 到 \(m_g\) 的谱的距离。由于后者与以 \(\mu\) 为测度的 \(g\) 的本质像重合，结果是 IV, p. 182 的引理 3 的推论。

最后，断言 \(c\) ) 由 \(b\) ) 及 \(\mathrm{PSp}_{\varepsilon}(u)\) 的定义（IV, p. 250, 定义 9）得出。

推论 3. --- 设 u 是希尔伯特空间 E 上的正规部分算子。设 \(f \in \mathcal{L}_{\mathrm{u}}(\mathrm{Sp}(u))\) 。对 E 中所有 x 与 y，\((x, y)\) 关于 \(f(u)\) 的谱测度是像测度 \(f(\mu)\) ，其中 \(\mu\) 是 \((x, y)\) 关于 u 的谱测度。

鉴于谱定理（IV, p. 266 的定理 1），可以假设 \(u\) 是乘法算子 \(m_g\) （作用于 \(\mathrm{L}^2 (\mathrm{X},\mu)\) ），其中 X 是局部紧拓扑空间，\(\mu\) 是 X 上的正测度，且 \(g\in \mathcal{C}(\mathrm{X})\) 。设 \(f_{1}\) 与 \(f_{2}\) 属于 \(\mathcal{L}^2 (\mathrm{X},\mu)\) ，其类 \(\widetilde{f}_1\) 与 \(\widetilde{f}_2\) 在 \(\mathrm{L}^2 (\mathrm{X},\mu)\) 中。由于 \(f(m_g) = m_{f\circ g}\) （命题 4, c)），\((\widetilde{f}_1,\widetilde{f}_2)\) 关于 \(f(m_g)\) 的谱测度是像测度 \((f\circ g)(\overline{f}_1f_2\cdot \mu)\) （IV, p. 269 的引理 4, b)）。这个测度等于像测度 \(f(g(\overline{f}_1f_2\cdot \mu))\) （INT, V, p. 72, § 6, 第 4 目, 命题 4, a)），由此即得断言（IV, p. 269 的引理 4, b)）。

推论 4. --- 设 \(g \in \mathcal{L}_{\mathrm{u}}(\mathrm{Sp}(f(u)))\) 。我们有 \(g(f(u)) = (g \circ f)(u)\) 。

我们有 \(g \circ f \in \mathcal{L}_{\mathrm{u}}(\mathrm{Sp}(u))\) （IV, p. 184 的引理 5）。对 E 中所有 \(x\) 与 \(y\) ，用 \(\mu_{x,y}^{\prime}\) 表示 \((x,y)\) 关于 \(f(u)\) 的谱测度。由前面的推论以及 INT, V, p. 71, § 6, 第 2 目, 定理 1，我们有 \(g \in \mathcal{L}^{2}(\mathrm{Sp}(f(u)), \mu_{x,y}^{\prime})\) 当且仅当 \(g \circ f \in \mathcal{L}^{2}(\mathrm{Sp}(u), \mu_{x,y})\) ，因此 \(g(f(u))\) 的定义域等于 \((g \circ f)(u)\) 的定义域。于是，对所有 \(x \in \mathrm{E}\) 及 \(y \in \operatorname{dom}(g(f(u)))\) ，有公式

\[
\begin{array}{c} \langle x \mid g (f (u)) y \rangle = \int_ {\mathrm{Sp} (f (u))} g \mu_ {x, y} ^ {\prime} \\ = \int_ {\mathrm{Sp} (f (u))} g f (\mu_ {x, y}) = \langle x \mid (g \circ f) y \rangle \end{array}
\]

（loc. cit.），由此 \(g(f(u)) = (g \circ f)(u)\) 。

命题 5. --- 设 u 是希尔伯特空间 E 上的正规部分算子。

\begin{enumerate}
\def\labelenumi{\alph{enumi})}
\item
  \(If f \in \mathcal{L}_{\mathrm{u}}^{\infty}(\mathrm{Sp}(u)),\) 则 \(f(u) \in \mathcal{L}(\mathrm{E})\) ；
\item
  \(If f \in \mathcal{L}_{\mathrm{u}}^{\infty}(\mathrm{Sp}(u))\) 且 \(g \in \mathcal{L}_{\mathrm{u}}(\mathrm{Sp}(u))\) ，则 \(f(u) \circ g(u) \subset (fg)(u)\) ；
\item
  映射 \(f \mapsto f(u)\) 从 \(\mathcal{L}_{\mathrm{u}}^{\infty}(\mathrm{Sp}(u))\) 到 \(\mathcal{L}(\mathrm{E})\) 中是星代数的连续单位态射。特别地，有 \(\|f(u)\|\leqslant\|f\|_{\infty}\) ，其中 \(f\in\mathcal{L}_{\mathrm{u}}^{\infty}(\mathrm{Sp}(u))\) ；
\item
  若 \(u \in \mathcal{L}(\mathrm{E})\) ，则对所有 \(f \in \mathcal{L}_{\mathrm{u}}^{\infty}(\mathrm{Sp}(u))\) ，自同态 \(f(u)\) 属于 \(u\) 在 \(\mathcal{L}(\mathrm{E})\) 中的双交换子。
\end{enumerate}

设 \(f \in \mathcal{L}_{\mathrm{u}}^{\infty}(\mathrm{Sp}(u))\) 。由 p. 272 的公式 (8)，我们有 \(\mathrm{D}_f = \mathrm{E}\) ，且 \(f(u)\) 是 \(\mathrm{E}\) 的连续自同态，由此得断言 \(a\) )。

设 \(g \in \mathcal{L}_{\mathrm{u}}(\mathrm{Sp}(u))\) 。设 \(y \in \mathrm{D}_g\) 。我们有 \(y \in \mathrm{D}_{fg}\) ，并且对所有 \(x \in \mathrm{E}\) ，由 p. 271 的公式 (7) 可得 \(\langle \overline{f}(u)x | g(u)y \rangle = \langle x | (fg)(u)y \rangle\) ，由此 \(f(u)(g(u)y) = (fg)(u)(y)\) ，这就证明了 \(b\) 。

由 \(a\) ) 与 \(b\) )，映射 \(f \mapsto f(u)\) 从 \(\mathcal{L}_{\mathrm{u}}^{\infty}(\mathrm{Sp}(u))\) 到 \(\mathcal{L}(\mathrm{E})\) 中，是对合代数的单位态射；因而它是范数 \(\leqslant 1\) 的连续态射（I, p. 104, 命题 2），由此得断言 \(c\) )。

假设 \(u\) 是 E 的自同态。设 \(v \in \mathcal{L}(\mathrm{E})\) 与 \(u\) 交换。于是，由连续函数演算的性质（I, p. 110, 注），我们有 \(v \circ f(u) = f(u) \circ v\) ，这里 \(f \in \mathcal{C}(\mathrm{Sp}(u))\) 。设 \(x\) 与 \(y\) 属于 E。公式

\[
\langle x \mid (v \circ f (u)) y \rangle = \langle x \mid (f (u) \circ v) y \rangle ,\tag{9}
\]

对每个函数 \(f \in \mathcal{C}(\mathrm{Sp}(u))\) 成立，它表示 \((v^{*}(x), y)\) 与 \((x, v(y))\) 关于 u 的谱测度相等。由 p. 271 的公式 (6)，这个等式蕴含公式 (9) 对所有 \(f \in \mathcal{L}_{\mathrm{u}}^{\infty}(\mathrm{Sp}(u))\) 成立。因此我们有 \(v \circ f(u) = f(u) \circ v\) 。

推论. --- 设 \(f\) 与 \(g\) 属于 \(\mathcal{L}_{\mathrm{u}}(\mathrm{Sp}(u))\) ，并设 \((x,y)\in\mathrm{D}_{f}\times\mathrm{D}_{g}\) 。\((f(u)x,g(u)y)\) 关于 \(u\) 的谱测度是测度 \(\overline{f}g\cdot\mu_{x,y}\) ，其中 \(\mu_{x,y}\) 是 \((x,y)\) 关于 \(u\) 的谱测度。

用 \(\nu\) 表示 \((f(u)x,g(u)y)\) 关于 \(u\) 的谱测度。对每个函数 \(\varphi \in \mathcal{H}(\mathrm{Sp}(u))\) ，我们有

\[
\begin{array}{r l} & {\int_ {\mathrm{Sp} (u)} \varphi \nu = \langle f (u) x | \varphi (u) (g (u) y) \rangle} \\ & {\qquad = \langle f (u) x | (\varphi g) (u) y \rangle = \int_ {\mathrm{Sp} (u)} \overline {{f}} \varphi g \mu_ {x, y},} \end{array}
\]

（命题 5, b)）由此 \(\nu = \overline{f} g \cdot \mu_{x,y}\) 。

命题 6. --- 设 u 是希尔伯特空间 E 上的正规部分算子。设 \((f_n)_{n \in \mathbf{N}}\) 是 \(\mathcal{L}_{\mathrm{u}}(\mathrm{Sp}(u))\) 中的序列，它逐点收敛于 \(f \in \mathcal{L}_{\mathrm{u}}(\mathrm{Sp}(u))\) ，且存在 \(g \in \mathcal{L}_{\mathrm{u}}(\mathrm{Sp}(u))\) 满足 \(|f_n| \leqslant g\) ，对所有 \(n \in \mathbf{N}\) 。这时我们有 \(\operatorname{dom}(g(u)) \subset \operatorname{dom}(f_n(u))\) ，对所有 \(n \in \mathbf{N}\) ，以及 \(\operatorname{dom}(g(u)) \subset \operatorname{dom}(f(u))\) 。此外，对 \(g(u)\) 的定义域的每个元素 y，我们有

\[
f (u) (y) = \lim _ {n \to + \infty} f _ {n} (u) (y).
\]

特别地，若 \(f_{n} \in \mathcal{L}_{\mathrm{u}}^{\infty}(\mathrm{Sp}(u))\) 对所有 \(n \in \mathbf{N}\) 成立，且函数 \(f_{n}\) 一致有界，则 \(f_{n}(u)\) 收敛于 \(f(u)\) ，这里的空间 \(\mathcal{L}(\mathrm{E})\) 赋以逐点收敛拓扑。

用 \(\mu_{x,y}\) （相应地，\(\mu_{x}\) ）表示 \((x,y)\) （相应地，x）关于 u 的谱测度。

设 \(y \in \mathrm{dom}(g(u))\) ，于是 \(g \in \mathcal{L}^2(\mathrm{Sp}(u), \mu_y)\) 。条件 \(|f_n| \leqslant g\) 蕴含 \(f_n \in \mathcal{L}^2(\mathrm{Sp}(u), \mu_y)\) ，因此 \(y \in \mathrm{dom}(f_n(u))\) 。

由于 \((f_{n})\) 逐点收敛于 f 且 \(|f_{n}| \leqslant g\) ，由勒贝格定理（INT, IV, p. 137, § 3, 第 7 目, 定理 6），我们有 \(f \in \mathcal{L}^{2}(\mathrm{Sp}(u), \mu_{y})\) ，因此 \(y \in \mathrm{dom}(f(u))\) 。此外，序列 \((f_{n})\) 在 \(\mathcal{L}^{2}(\mathrm{Sp}(u), \mu_{y})\) 中收敛于 f，于是 \(f_{n}(u)y\) 的范数收敛于 \(f(u)y\) 的范数。

设 \(x \in \mathrm{E}\) 。函数 \(f\) 与 \(g\) ，以及函数 \(f_n\) （对所有 \(n \in \mathbf{N}\) ），都属于 \(\mathcal{L}^1(\mathrm{Sp}(u), \mu_{x,y})\) （命题 3）。由勒贝格定理（INT, IV, loc. cit.），序列 \((f_n)\) 收敛于 \(f\) （在 \(\mathcal{L}^1(\mathrm{Sp}(u), \mu_{x,y})\) 中），由此

\[
\langle x \mid f _ {n} (u) y \rangle = \int_ {\operatorname{Sp} (u)} f _ {n} \mu_ {x, y} \rightarrow \int_ {\operatorname{Sp} (u)} f \mu_ {x, y} = \langle x \mid f (u) y \rangle .
\]

由此得出 \(f_{n}(u)(y)\) 收敛于 \(f(u)(y)\) （EVT, V, p. 17, 命题 10）。

命题 7. --- 设 X 是局部紧拓扑空间，ν 是 X 上的测度。设 g: C × X → C 是具有紧支撑的连续函数。对 z ∈ C，置

\[
h (z) = \int_ {\mathrm{X}} g (z, x) d \nu (x).
\]

\begin{enumerate}
\def\labelenumi{\alph{enumi})}
\item
  从 C 到 C 的映射 h 是连续且有界的；
\item
  从 X 到 \(\mathcal{L}(\mathrm{E})\) 的、由 \(x\mapsto g(u,x)\) 定义的映射是 \(\nu\)-可积的，并且我们有
\end{enumerate}

\[
h (u) = \int_ {\mathrm{X}} g (u, x) d \nu (x).
\]

函数 h 是有界的，因为 g 是具有紧支撑的连续函数；且由 INT, IV, p. 144, § 4, 第 3 目, 推论 1，它是连续的。由于 g 连续且具有紧支撑，映射 \(x \mapsto g(u, x)\) 从 X 到 \(\mathcal{L}(\mathrm{E})\) 中连续。这个映射具有紧支撑，故在 X 上关于 \(\nu\) 有界且可积。我们记

\[
v = \int_ {\mathrm{X}} g (u, x) d \nu (x) \in \mathscr {L} (\mathrm{E}).
\]

设 y 与 z 是 E 的元素，并设 \(\mu\) 是 \((y, z)\) 关于 u 的谱测度。我们有

\[
\langle y \mid v (z) \rangle = \int_ {\mathrm{X}} \langle y \mid g (u, x) z \rangle d \nu (x) = \int_ {\mathrm{X}} \left(\int_ {\operatorname{Sp} (u)} g (\lambda , x) d \mu (\lambda)\right) d \nu (x)
\]

（p. 271 的公式 (6)）。由于 \(g \in \mathcal{K}(\mathbf{C} \times \mathbf{X})\) ，由此可得

\[
\langle y \mid v (z) \rangle = \int_ {\operatorname{Sp} (u)} \Bigl (\int_ {\mathrm{X}} g (\lambda , x) d \nu (x) \Bigr) d \mu (\lambda) = \langle y \mid h (u) z \rangle
\]

依据 INT, III, p. 84, § 4, 第 1 目, 定理 2 以及 p. 271 的公式 (6)。这就证明了 \(v = h(u)\) ，此即所要证者。

\subsection{谱投影}

设 \(u\) 是复希尔伯特空间 E 上的正规部分算子。设 A 是 \(\mathrm{Sp}(u)\) 的普遍可测子集，\(\varphi_{\mathrm{A}}\) 是它的特征函数。由于 \(\varphi_{\mathrm{A}}\) 有界且满足 \(\varphi_{\mathrm{A}}^{2} = \varphi_{\mathrm{A}}\) ，E 的自同态 \(\varphi_{\mathrm{A}}(u)\) 是 E 的正交投影。它称为由 A 定义的 \(u\) 的谱投影。我们记 \(p_{\mathrm{A}} = \varphi_{\mathrm{A}}(u)\) 。我们有 \(p_{\emptyset} = 0\) 及 \(p_{\mathrm{Sp}(u)} = 1_{\mathrm{E}}\) 。

设 A 是 C 的普遍可测子集。由 A 定义的 u 的谱投影是谱投影 \(p_{\mathrm{Sp}(u)\cap\mathrm{A}}\) 。它也简单地记作 \(p_{A}\) 。对每个函数 \(f \in \mathcal{L}_{\mathrm{u}}(\mathrm{Sp}(u))\) ，对所有 \(x \in E\) 及每个 \(y \in \mathrm{dom}(f(u))\) ，我们有公式

\[
\langle p _ {\mathrm{A}} (x) \mid f (u) y \rangle = \int_ {\operatorname{Sp} (u) \cap \mathrm{A}} f d \mu ,\tag{10}
\]

其中 \(\mu\) 是 \((x,y)\) 关于 \(u\) 的谱测度（p. 271 的公式 (7)）。设 A 与 B 是 \(\mathrm{Sp}(u)\) 的普遍可测子集。由于 \(\varphi_{\mathrm{A}}\varphi_{\mathrm{B}} = \varphi_{\mathrm{A}\cap \mathrm{B}}\) ，我们有 \(p_{\mathrm{A}} \circ p_{\mathrm{B}} = p_{\mathrm{A}\cap \mathrm{B}}\) （IV, p. 275 的命题 5, \(c\) )）。特别地，若 A 与 B 不相交，则 \(p_{\mathrm{A}}\) 的像与 \(p_{\mathrm{B}}\) 的像正交。

命题 8. --- 设 \((\mathrm{A}_{i})_{i\in\mathrm{I}}\) 是 \(\mathrm{Sp}(u)\) 的两两不相交的普遍可测子集组成的可数族，并设 \(p_{i}\) 是由 \(A_{i}\) 定义的 u 的谱投影。诸集合 \(A_{i}\) 的并 A 是 \(\mathrm{Sp}(u)\) 的普遍可测子集，且级数 \(\sum_{i}p_{i}\) 收敛于 \(p_{A}\) ，这里的 \(\mathcal{L}(\mathrm{E})\) 赋以逐点收敛拓扑。

由 INT, IV, p. 177, § 5, 第 4 目, 推论 2，集合 A 是普遍可测的。级数 \(\sum_{i} \varphi_{\mathrm{A}_{i}}\) 逐点收敛于 \(\varphi_{\mathrm{A}}\) ，且其部分和以 1 为界。因此断言由 IV, p. 276 的命题 6 得出。

命题 9. --- 设 A 是 Sp(u) 的闭子集。设 \(\varphi_{\mathrm{A}}\) 是 A 的特征函数，\(p_{\mathrm{A}} = \varphi_{\mathrm{A}}(u)\) 是由 A 定义的 u 的谱投影。用 \(\mathrm{E}_{\mathrm{A}}\) 表示 \(p_{\mathrm{A}}\) 的像。它是 E 的闭子空间。

\begin{enumerate}
\def\labelenumi{\alph{enumi})}
\item
  子空间 \(E_{A}\) 是使 x 关于 u 的谱测度的支撑包含在 A 中的那些 \(x \in E\) 所成的集；
\item
  若 A 在 \(\mathbf C\) 中有界，则 \(E_{A}\) 包含在 u 的定义域中；
\item
  对属于 u 的定义域的所有 x，我们有 \(p_{\mathrm{A}}(x) \in \operatorname{dom}(u)\) 且 \(u(p_{\mathrm{A}}(x)) \in \mathrm{E}_{\mathrm{A}}\) ；特别地，\(u(x) \in \mathrm{E}_{\mathrm{A}}\) ，只要 \(x \in \operatorname{dom}(u) \cap \mathrm{E}_{\mathrm{A}}\) 。
\end{enumerate}

对 E 中的 \(x\) ，我们用 \(\mu_x\) 表示 \(x\) 关于 \(u\) 的谱测度。

我们来证明 \(a\) )。设 \(x\) 是 \(\mathrm{E}_{\mathrm{A}}\) 的元素。设 \(z \in \mathbf{C} - \mathbf{A}\) ，并设 U 是 \(z\) 的与 A 不相交的相对紧开邻域。对每个支撑包含在 U 中的函数 \(f \in \mathcal{K}(\mathbf{C})\) ，我们有

\[
\int_ {\mathrm{Sp} (u)} f \mu_ {x} = \langle x | f (u) x \rangle = \langle p _ {\mathrm{A}} (x) | f (u) x \rangle = \int_ {\mathrm{Sp} (u) \cap \mathrm{A}} f \mu_ {x} = 0
\]

（依据公式 (10)）。这表明 \(z\) 不属于 \(\mu_x\) 的支撑。因此，\(\mu_x\) 的支撑包含在 A 中。

反之，设 \(x \in E\) 使 \(\mu_{x}\) 的支撑包含在 A 中。由此可得

\[
\langle x \mid p _ {\mathrm{A}} (x) \rangle = \int_ {\operatorname{Sp} (u) \cap \mathrm{A}} \mu_ {x} = \int_ {\operatorname{Sp} (u)} \mu_ {x} = \langle x \mid x \rangle ,
\]

\((loc. \ cit.) \ \text{因此} \ \|p_{\mathrm{A}}(x)-x\|^{2} = \|p_{\mathrm{A}}(x)\|^{2} - \|x\|^{2} \leqslant 0 \ \text{从而 } p_{\mathrm{A}}(x) = x, \ \text{即 } x \in \mathrm{E}_{\mathrm{A}}.\)

断言 b) 由 a) 及 IV, p. 273 的注得出。

我们来证明 \(c\) )。\(u\) 的定义域是这样的 \(x \in \mathrm{E}\) 所成的集：\(\operatorname{Sp}(u)\) 的恒等函数属于 \(\mathcal{L}^2(\operatorname{Sp}(u), \mu_x)\) （loc. cit.）。由于 \(\mu_{p_{\mathrm{A}}(x)} = \varphi_{\mathrm{A}} \cdot \mu_x\) 对 \(x \in \mathrm{E}\) 成立（IV, p. 275 的命题 5 的推论），我们有 \(p_{\mathrm{A}}(x) \in \operatorname{dom}(u)\) ，只要 \(x \in \operatorname{dom}(u)\) 。

设 \(x \in \operatorname{dom}(u)\) 且 \(y = p_{\operatorname{A}}(x)\) ；我们有 \(y \in \operatorname{dom}(u) \cap \operatorname{E}_{\operatorname{A}}\) ，依据 \(c\) )。\(u(y)\) 关于 \(u\) 的谱测度是 \(|\operatorname{Id}_{\operatorname{Sp}(u)}|^2 \cdot \mu_y\) （loc. cit.）。由于 \(\mu_y\) 的支撑在 A 中，\(\mu_{u(y)}\) 亦然，故 \(u(y) \in \operatorname{E}_{\operatorname{A}}\) ，依据 \(a\) )。

推论. --- 设 \(\lambda \in \mathrm{Sp}(u)\) 。\(p_{\{\lambda\}}\) 的像是 \(u\) 关于 \(\lambda\) 的特征子空间。

设 \(x \in \operatorname{dom}(u)\) 。用 \(\mu\) （相应地，\(\nu\) ）表示元素 \(x\) 关于 \(u\) 的谱测度（相应地，\((x, u(x))\) 关于 \(u\) 的谱测度）。我们有 \(\nu = \operatorname{Id}_{\operatorname{Sp}(u)} \cdot \mu\) （IV, p. 275 的命题 5 的推论）。若 \(u(x) = \lambda x\) ，我们还有 \(\nu = \lambda \mu\) ，由此得等式 \(\operatorname{Id}_{\operatorname{Sp}(u)} \cdot \mu = \lambda \mu\) 。这蕴含 \(\mu\) 的支撑包含在 \(\{\lambda\}\) 中（参见 INT, V, p. 46, § 5, 第 3 目, 推论 2），因此 \(x\) 属于 \(p_{\{\lambda\}}\) 的像（命题 9, a)）。

反之，假设 \(x\) 属于 \(\mathrm{E}_{\lambda}\) ，即 \(p_{\{\lambda\}}\) 的像。那么 \(x\) 属于 \(\operatorname{dom}(u)\) ，且 \(u(x)\) 也属于 \(\mathrm{E}_{\lambda}\) （loc. cit., \(b\) )）。由于 \(\varphi_{\{\lambda\}} \cdot (\operatorname{Id}_{\operatorname{Sp}(u)} - \lambda) = 0\) ，我们有关系式 \(p_{\{\lambda\}} \circ (u - \lambda 1_{\mathrm{E}}) \subset 0\) （IV, p. 275 的命题 5, \(c\) )），由此 \(0 = p_{\{\lambda\}}(u(x)) - \lambda p_{\{\lambda\}}(x) = u(x) - \lambda x\) 。

命题 10. — 设 \(\lambda \in \mathbf{C}\) 。我们有 \(\lambda \in \mathrm{Sp}(u)\) 当且仅当：对 \(\lambda\) 在 \(\mathbf{C}\) 中的每个开邻域 V，u 关于 V 的谱投影 \(p_{\mathrm{V}}\) 非零。

若 \(\lambda \notin \mathrm{Sp}(u)\) ，则存在开邻域 \(V\) ，它包含 \(\lambda\) ，位于 \(\mathbf{C}\) 中，且与 \(\mathrm{Sp}(u)\) 不相交，于是 \(p_{\mathrm{V}} = p_{\varnothing} = 0\) 。

反之，假设存在 \(\lambda\) 在 \(\mathbf{C}\) 中的开邻域 V 使得 \(p_{V} = 0\) 。设 c > 0 使得以 \(\lambda\) 为中心、以 c 为半径的圆盘包含在 V 中。

设 \(x \in \mathrm{dom}(u)\) ，并设 \(\mu_x\) 是 \(x\) 关于 \(u\) 的谱测度。由于 \(\mu_x(\mathrm{V}) = \langle x | p_{\mathrm{V}}(x) \rangle = 0\) ，我们计算出

\[
\begin{array}{r l} & {\int_ {\mathbf {C}} | z - \lambda | ^ {2} d \mu_ {x} (z) = \int_ {\mathbf {C - V}} | z - \lambda | ^ {2} d \mu_ {x} (z)} \\ & {\qquad \geqslant c ^ {2} \int_ {\mathbf {C - V}} d \mu_ {x} (z) = c ^ {2} \int_ {\mathbf {C}} d \mu_ {x} (z) = c ^ {2} \| x \| ^ {2}.} \end{array}
\]

但是，另一方面，我们有

\[
\| u (x) - \lambda x \| ^ {2} = \int_ {\mathbf {C}} | z - \lambda | ^ {2} d \mu_ {x} (z) = \| u ^ {*} (x) - \overline {{\lambda}} x \| ^ {2}
\]

（p. 272 的公式 (8)），由此 \(\|u(x)-\lambda x\|\geqslant c\|x\|\) 且 \(\|u^{*}(x)-\overline{\lambda}x\|\geqslant c\|x\|\) 。于是 \(\lambda\) 属于 u 的预解集（IV, p. 248 的引理 5）。证明完毕。

推论. — 设 A 是 \(\mathbf{C}\) 中的开集，且 \(n \in \mathbf N\) 。若 A 含有 \(\mathrm{Sp}(u)\) 的 n 个元素，则 u 关于 A 的谱投影 \(p_{A}\) 的像的维数至少等于 n。

设 \(\lambda_1, \ldots, \lambda_n\) 是 \(u\) 的谱中属于 A 的互不相同的元素。存在两两不相交的开集族 \((\mathrm{V}_i)_{1 \leqslant i \leqslant n}\) （\(\mathbf{C}\) 中的开集），使得 \(\lambda_i \in \mathrm{V}_i\) 对 \(1 \leqslant i \leqslant n\) 成立。设 B 是诸集合 \(\mathrm{V}_i\) 的并。\(p_{\mathrm{A}}\) 的像包含谱投影 \(p_{\mathrm{B}}\) 的像；此外，由于 \(p_{\mathrm{B}}\) 是诸投影 \(p_{\mathrm{V}_i}\) 之和，且 \(p_{\mathrm{V}_i}\) 的像与 \(p_{\mathrm{V}_j}\) 的像正交（对所有 \(i \neq j\) ），结果由命题 10 得出。

\subsection{埃尔弗--舍斯特兰德公式}

在本节中，E 表示一个复希尔伯特空间。我们将对自伴部分算子的函数演算的某些情形得到一个公式，它直接用所论算子的预解式来表达。

我们给 R （相应地，C）赋以勒贝格测度，记作 \(\mu\) ，并用映射 \((x,y)\mapsto\exp(2i\pi xy)\) 把群 R 与它的对偶群等同起来（参见 II, p. 236 的推论 3）。

对函数 \(f\) ——它在 \(\mathbf{R}^{2}\) 的开集 U 上定义且可微，该开集等同于 \(\mathbf{C}\) ，赋以实坐标 \(x\) 与 \(y\) ——我们记

\[
{\frac {\partial f}{\partial \overline {{{z}}}}} = {\frac {1}{2}} {\Bigl (} {\frac {\partial f}{\partial x}} + i {\frac {\partial f}{\partial y}} {\Bigr)}
\]

（参见 VAR, R2, 8.8.10, p. 24）。

引理 5. --- 设 \(f \in \mathcal{D}(\mathbf{R})\) 。存在函数 \(\tilde{f}\) 属于 \(\mathcal{D}(\mathbf{C})\) ，它与 \(f\) 在 \(\mathbf{R}\) 上重合，且满足

\[
\frac {\partial \widetilde {f}}{\partial \overline {{z}}} (x, 0) = 0\tag{11}
\]

对所有 \(x \in \mathbf{R}\) 成立。这时我们有

\[
\frac {\partial \widetilde {f}}{\partial y} (x, 0) = i f ^ {\prime} (x)\tag{12}
\]

对所有 \(x \in R\) 成立，且存在实数 C ≥ 0 使得

\[
\left| \frac {\partial \widetilde {f}}{\partial \overline {{z}}} (x, y) \right| \leqslant C | y |\tag{13}
\]

对所有 \((x,y)\in \mathbf{R}^2\)

存在 \(\varphi\in\mathcal{D}(\mathbf{R})\) ，其支撑包含在 \([-2,2]\) 中，且在 \([-1,1]\) 上等于 1（IV, p. 196 的引理 1）。置

\[
\widetilde {f} (x, y) = \bigl (f (x) + i y f ^ {\prime} (x) \bigr) \varphi (y)
\]

对 \((x,y)\in\mathbf{R}^{2}\) 。我们有 \(\widetilde{f}\in\mathcal{D}(\mathbf{C})\) ，且 \(\widetilde{f}\) 在 R 上与 f 重合。此外，对任意 \((x,y)\in\mathbf{R}^{2}\) ，可得

\[
\frac {\partial \widetilde {f}}{\partial \overline {{z}}} (x, y) = \frac {1}{2} \Big ((i f (x) - y f ^ {\prime} (x)) \varphi^ {\prime} (y) + i y f ^ {\prime \prime} (x) \varphi (y) \Big).
\]

由于函数 \(\varphi\) 在 0 的一个邻域中等于 1，我们有 \(\varphi'(0) = 0\) ，由此得 (11)。

设 \(\tilde{f}\) 是 \(\mathcal{D}(\mathbf{C})\) 中满足 (11) 的函数。由此即得公式 (12)，而估计式 (13) 则借助中值定理（FVR, I, p. 23, 定理 2）得到。

我们称 \(\tilde{f}\) 是 f 的一个几乎解析延拓。

引理 6. --- 设 \(\varepsilon > 0\) 。函数 \(\sigma_{\varepsilon}\) 定义在 \(\mathbf{R}\) 上，由下式给出

\[
\sigma_ {\varepsilon} (x) = \frac {2 i \varepsilon x}{x ^ {2} + \varepsilon^ {2}}
\]

属于 \(L^{2}(\mathbf{R})\) 。它的傅里叶变换是 \(L^{2}(\mathbf{R})\) 中函数 \(\eta_{\varepsilon}\) 的类，该函数在 0 处取值 0 且满足

\[
\eta_ {\varepsilon} (y) = \frac {2 \pi \varepsilon y}{| y |} e ^ {- 2 \pi \varepsilon | y |}
\]

对所有 \(y \neq 0\) 成立。

由 IV, p. 199 的命题 3，有 \(\sigma_{\varepsilon} \in \mathrm{L}^{2}(\mathbf{R})\) ，且函数 \(\eta_{\varepsilon}\) 的类属于 \(\mathrm{L}^{2}(\mathbf{R}) \cap \mathrm{L}^{1}(\mathbf{R})\) 。对每个 \(x \in \mathbf{R}\) ，我们有

\[
\begin{array}{r} \overline {{\mathcal {F}}} (\eta_ {\varepsilon}) (x) = 2 \pi \varepsilon \int_ {\mathbf {R} _ {+}} e ^ {2 \pi (i x - \varepsilon) y} d y - 2 \pi \varepsilon \int_ {\mathbf {R} _ {-}} e ^ {2 \pi (i x + \varepsilon) y} d y \\ = \frac {\varepsilon}{\varepsilon - i x} - \frac {\varepsilon}{\varepsilon + i x} = \frac {2 i \varepsilon x}{x ^ {2} + \varepsilon^ {2}}, \end{array}
\]

由此 \(\overline{\mathcal{F}}(\eta_{\varepsilon}) = \sigma_{\varepsilon}\) 。于是结果由 \(\mathrm{L}^2(\mathbf{R})\) 中的傅里叶反演公式（II, p. 220 的定理 2 的推论）得出。

引理 7. --- 设 \(f \in \mathcal{D}(\mathbf{R})\) ，并设 \(\widetilde{f} \in \mathcal{D}(\mathbf{C})\) 是 \(f\) 的一个几乎解析延拓。对 \(\varepsilon > 0\) ，定义 \(f_{\varepsilon}\) 于 \(\mathbf{R}\) 上如下

\[
f _ {\varepsilon} (x) = - \frac {1}{2 i \pi} \int_ {\mathbf {R}} \biggl (\frac {\widetilde {f} (y + i \varepsilon)}{y - x + i \varepsilon} - \frac {\widetilde {f} (y - i \varepsilon)}{y - x - i \varepsilon} \biggr) d y.
\]

那么 \(f_{\varepsilon}\) 是连续且有界的，且 \(f_{\varepsilon}\) 收敛于 \(f\) （在 \(\mathcal{C}_b(\mathbf{R})\) 中，当 \(\varepsilon\) 趋于 0 时）。

\(f_{\varepsilon}\) 的连续性是 INT, IV, p. 144, § 4, 第 3 目, 推论 1 的推论，因为 \(\widetilde{f}\) 具有紧支撑。此外，若 \(r\) 使得 \(\widetilde{f}\) 的支撑包含在 \([-r, r] \times \mathbf{R}\) 中，我们有

\[
| f _ {\varepsilon} (x) | \leqslant 2 \| \widetilde {f} \| _ {\infty} \int_ {- r} ^ {r} \frac {1}{\sqrt {(x - y) ^ {2} + \varepsilon^ {2}}} d y \leqslant \frac {4 r}{\varepsilon} \| \widetilde {f} \| _ {\infty},
\]

因此 \(f_{\varepsilon}\) 是有界的。

一阶泰勒展开（FVR, I, p. 30, 命题 3）与公式 (13) 表明，存在 M ≥ 0 及一个函数 \(\varrho_{1}\) ，它取值于 \(R^{2}\) 中，使得

\[
\widetilde {f} (y + i \gamma) = f (y) + i \gamma f ^ {\prime} (y) + \gamma^ {2} \varrho_ {1} (y; \gamma), \text {   and   } | \varrho_ {1} (y; \gamma) | \leqslant M,
\]

对所有 \(y \in \mathbf{R}\) 及 \(\gamma \in \mathbf{R}\) 成立。由于映射 \(y \mapsto \widetilde{f}(y + i\gamma)\) 的支撑包含在 \([-r, r]\) 中（对所有 \(\gamma \in \mathbf{R}\) ），映射 \(y \mapsto \varrho_{1}(y; \gamma)\) 的支撑也包含在 \([-r, r]\) 中，对所有 \(\gamma \in \mathbf{R}\) 。

设 \(g_{\varepsilon}\) 是定义在 \(\mathbf{R}^2\) 上的函数

\[
g _ {\varepsilon} (x, y) = \frac {\widetilde {f} (y + i \varepsilon)}{y - x + i \varepsilon} - \frac {\widetilde {f} (y - i \varepsilon)}{y - x - i \varepsilon}.
\]

对每个 \((x,y)\in \mathbf{R}^2\) 及 \(\varepsilon >0\) ，我们得到

\[
\begin{array}{r l} g _ {\varepsilon} (x, y) = - \frac {2 i \varepsilon}{(x - y) ^ {2} + \varepsilon^ {2}} f (y) + \frac {2 i \varepsilon (y - x)}{(x - y) ^ {2} + \varepsilon^ {2}} f ^ {\prime} (y) \\ & \quad + \varepsilon^ {2} \bigg (\frac {\varrho_ {1} (y ; \varepsilon)}{y - x + i \varepsilon} - \frac {\varrho_ {1} (y ; - \varepsilon)}{y - x - i \varepsilon} \bigg). \end{array}
\]

设 \(x \in \mathbf{R}\) 且 \(\varepsilon > 0\) 。我们有

\[
\begin{array}{c} \left| \varepsilon^ {2} \int_ {\mathbf {R}} \left(\frac {\varrho_ {1} (y ; \varepsilon)}{y - x + i \varepsilon} - \frac {\varrho_ {1} (y ; - \varepsilon)}{y - x - i \varepsilon}\right) d y \right| \leqslant 2 \mathrm{M} \varepsilon^ {2} \int_ {- r} ^ {r} \frac {d y}{\sqrt {(x - y) ^ {2} + \varepsilon^ {2}}} \\ \leqslant 4 \mathrm{Mr} \varepsilon . \end{array}
\]

因此，对每个 \(x \in \mathbf{R}\) ，可得

\[
f _ {\varepsilon} (x) = - \frac {1}{2 i \pi} \int_ {\mathbf {R}} g _ {\varepsilon} (x, y) d y = (f * \delta_ {\varepsilon}) (x) - \frac {1}{2 i \pi} (f ^ {\prime} * \sigma_ {\varepsilon}) (x) + k _ {\varepsilon} (x)
\]

其中 \(\delta_{\varepsilon}\) 与 \(\sigma_{\varepsilon}\) 是 \(\mathbf{R}\) 上由下式定义的函数

\[
\delta_ {\varepsilon} (x) = \frac {1}{\pi} \frac {\varepsilon}{x ^ {2} + \varepsilon^ {2}}, \qquad \sigma_ {\varepsilon} (x) = \frac {2 i \varepsilon x}{x ^ {2} + \varepsilon^ {2}},
\]

且 \(\| k_{\varepsilon}\|_{\infty}\leqslant 2\mathrm{Mr}\varepsilon .\)

函数 \(\mathcal{F}(f') \in \mathcal{S}(\mathbf{R})\) 是可积的（IV, p. 219 的命题 18 及 IV, p. 213 的命题 13）。由引理 6，可得

\[
\begin{array}{r l r} & & {\int_ {\mathbf {R}} | \mathcal {F} (f ^ {\prime}) (y) \mathcal {F} (\sigma_ {\varepsilon}) (y) | d y = 2 \pi \varepsilon \int_ {\mathbf {R}} | \mathcal {F} (f ^ {\prime}) (y) | e ^ {- 2 \pi \varepsilon | y |} d y} \\ & & {\leqslant 2 \pi \varepsilon \int_ {\mathbf {R}} | \mathcal {F} (f ^ {\prime}) (y) | d y,} \end{array}
\]

因此 \(\mathcal{F}(f')\mathcal{F}(\sigma_{\varepsilon})\) 在 \(\mathrm{L}^{1}(\mathbf{R})\) 中收敛于 0，当 \(\varepsilon\) 趋于 0 时。由于 \(f'\) 与 \(\sigma_{\varepsilon}\) 属于 \(\mathrm{L}^{2}(\mathbf{R})\) ，II, p. 223 的命题 14 蕴含 \(f'*\sigma_{\varepsilon}=\overline{\mathcal{F}}(\mathcal{F}(f')\mathcal{F}(\sigma_{\varepsilon}))\) 在 \(\mathcal{C}_{b}(\mathbf{R})\) 中收敛于 0。

对每个 \(\varepsilon > 0\) ，正测度 \(\delta_{\varepsilon} \cdot dx\) 在 \(\mathbf{R}\) 上的总质量为 1（参见 FVR, III, p. 7）。测度 \(\delta_{\varepsilon} \cdot dx\) （对 \(\varepsilon > 0\) ）所成的集，与 \(\mathbf{R}_{+}^{*}\) 中 0 的邻域滤子在该集上诱导的滤子，满足 INT, VIII, p. 137, § 2, 第 7 目的引理 4 的假设。因此，测度 \(\delta_{\varepsilon} \cdot dx\) 在 \(\mathcal{M}^{1}(\mathbf{R})\) 中收敛于点测度 \(\varepsilon_{0}\) ，当 \(\varepsilon\) 趋于 0 时。由此可得 \(f * \delta_{\varepsilon} \to f\) 在 \(\mathcal{C}_{b}(\mathbf{R})\) 中成立（INT, VIII, p. 163, § 4, 第 4 目）。引理得证。

我们用 \(\mu\) 表示 \(\mathbf{C}\) 上的勒贝格测度。

定理 2（埃尔弗--舍斯特兰德）. --- 设 u 是 E 上的自伴部分算子。设 \(f \in \mathcal{D}(\mathbf{R})\) ，且 \(\widetilde{f} \in \mathcal{D}(\mathbf{C})\) 是 \(f\) 的一个几乎解析延拓。设 h 是由下式定义的从 \(\mathbf{C}\) 到 \(\mathcal{L}(\mathrm{E})\) 中的映射

\[
h (\lambda) = \frac {\partial \widetilde {f}}{\partial \overline {{z}}} (\lambda) \mathrm{R} (u, \lambda)
\]

若 \(\lambda\in\mathbf{C} - \mathbf{R}\) ；而 \(h(\lambda)=0\) ，若 \(\lambda\in\mathbf{R}\) 。那么 h 关于 \(\mu\) 在 \(\mathbf{C}\) 上可积，且

\[
f (u) = - \frac {1}{\pi} \int_ {\mathbf {C}} h (\lambda) d \mu (\lambda).
\]

映射 h 是可测的，且其支撑是紧的。由于 u 是自伴的，我们有 \(\|\mathrm{R}(u,\lambda)\| \leqslant |\mathcal{I}(\lambda)|^{-1}\) ，对所有 \(\lambda \in \mathbf{C} - \mathbf{R}\) （IV, p. 248 的命题 17）。因此，依据公式 (13)，映射 h 是有界的，从而它在 \(\mathbf{C}\) 上可积。

设 \(\varepsilon \in \mathbf{R}_{+}^{*}\) 。我们用 \(\mathrm{F}_{\varepsilon}^{+}\) （相应地，\(\mathrm{F}_{\varepsilon}^{-}\) ）表示 \(\mathbf{C}\) 中由这样的 \(\lambda \in \mathbf{C}\) 所成的闭集：\(\mathcal{I}(\lambda) \geqslant \varepsilon\) （相应地，\(\mathcal{I}(\lambda) \leqslant -\varepsilon\) ）。我们有

\[
\int_ {\mathbf {C}} h (\lambda) d \mu (\lambda) = \lim _ {\varepsilon \to 0} \Bigl (\int_ {\mathrm{F} _ {\varepsilon} ^ {+}} h (\lambda) d \mu (\lambda) + \int_ {\mathrm{F} _ {\varepsilon} ^ {-}} h (\lambda) d \mu (\lambda) \Bigr).
\]

设 r > 0 使得 h 的支撑包含在 \(\mathbf{C} = [-r, r]^{2}\) 中。用 \(R_{\varepsilon}^{+}\) 表示 \([-2r, 2r] \times [\varepsilon, 2r]\) 这一 \(\mathbf{C}\) 中的长方体。它是 \(\mathbf{C}\) 的局部多面体子集（VAR, R2, 11.3, p. 48）。它满足下列条件：

\begin{enumerate}
\def\labelenumi{(\roman{enumi})}
\item
  我们有 \(R_{\varepsilon}^{+} \subset 2C \cap F_{\varepsilon}^{+}\) ；
\item
  集合 \(R_{\varepsilon}^{+}\) 包含 \(F_{\varepsilon}^{+}\) 与 h 的支撑的交；
\item
  正则边界 \(\partial R_{\varepsilon}^{+}\) （VAR, R2, 11.3.2, p. 49）包含线段 \(S_{\varepsilon} = [-r, r] + i\varepsilon \subset C\) ；
\item
  有 \(h(\lambda) = 0\) ，只要 \(\lambda \in \partial \mathrm{R}_{\varepsilon}^{+} - \mathrm{S}_{\varepsilon}\) 。
\end{enumerate}

用 \(d\lambda\) （相应地，\(d\overline{\lambda}\) ）表示 \(\mathbf{C}\) 上的、\(\mathbf{C}\) 的恒等映射（相应地，复共轭）的微分 1-形式。用 \(g\) 表示 \(\mathbf{C} - \mathbf{R}\) 上取值于 \(\mathcal{L}(\mathrm{E})\) 的、使得 \(g(\lambda) = \widetilde{f}(\lambda)\mathrm{R}(u,\lambda)\) （对 \(\lambda \in \mathbf{C} - \mathbf{R}\) ）的函数。设 \(\omega = g d\lambda\) ；它是 \(\mathbf{C} - \mathbf{R}\) 上具有紧支撑且取值于 \(\mathcal{L}(\mathrm{E})\) 的 1 次微分形式。由于 \(u\) 的预解式是全纯的（IV, p. 246 的命题 14），我们有

\[
d \omega = \left(\frac {\partial \widetilde {f}}{\partial \overline {{z}}} (\lambda) \mathrm{R} (u, \lambda) + \widetilde {f} (\lambda) \frac {\partial}{\partial \overline {{z}}} \mathrm{R} (u, \lambda)\right) d \overline {{\lambda}} \wedge d \lambda = - h (\lambda) d \lambda \wedge d \overline {{\lambda}}.
\]

与 \(d\omega\) 相伴的向量测度（VAR, R2, 10.4.3, p. 43）是关于勒贝格测度以 -2ih 为密度的测度。对局部多面体集 \(R_{\varepsilon}^{+}\) 及微分形式 \(\omega\) 应用斯托克斯公式（VAR, R2, 11.3.4, p. 49），我们便得到

\[
\frac {i}{2} \int_ {\mathrm{R} _ {\varepsilon} ^ {+}} d \omega = \frac {i}{2} \int_ {\partial \mathrm{R} _ {\varepsilon} ^ {+}} \omega = \frac {i}{2} \int_ {\mathrm{S} _ {\varepsilon}} \omega = \frac {i}{2} \int_ {- r} ^ {r} \widetilde {f} (y + i \varepsilon) \mathrm{R} (u, y + i \varepsilon) d y,
\]

由此

\[
\int_ {\mathrm{F} _ {\varepsilon} ^ {+}} h (\lambda) d \mu (\lambda) = - \frac {i}{2} \int_ {\mathbf {R}} \widetilde {f} (y + i \varepsilon) \mathrm{R} (u, y + i \varepsilon) d y.
\]

对 \(F_{\varepsilon}^{-}\) 作同样的推理，我们得到

\[
\int_ {\mathrm{F} _ {\varepsilon} ^ {-}} h (\lambda) d \mu (\lambda) = - \frac {i}{2} \int_ {\mathbf {R}} \widetilde {f} (y - i \varepsilon) \mathrm{R} (u, y - i \varepsilon) d y,
\]

并且我们得出，\(h\) 在 \(\mathbf{C}\) 上的积分是当 \(\varepsilon\to0\) 时下列量的极限

\[
v _ {\varepsilon} = \frac {i}{2} \int_ {\mathbf {R}} \Bigl (\widetilde {f} (y + i \varepsilon) \mathrm{R} (u, y + i \varepsilon) - \widetilde {f} (y - i \varepsilon) \mathrm{R} (u, y - i \varepsilon) \Bigr) d y.
\]

由命题 7，我们有 \(v_{\varepsilon} = \pi f_{\varepsilon}(u)\) ，其中 \(f_{\varepsilon}\) 是定义在 \(\mathbf{R}\) 上的函数

\[
f _ {\varepsilon} (x) = - \frac {1}{2 i \pi} \int_ {\mathbf {R}} \biggl (\frac {\widetilde {f} (y + i \varepsilon)}{y - x + i \varepsilon} - \frac {\widetilde {f} (y - i \varepsilon)}{y - x - i \varepsilon} \biggr) d y.
\]

由于 \(f_{\varepsilon} \to f\) 当 \(\varepsilon \to 0\) 时在 \(\mathcal{C}_{b}(\mathbf{R})\) 中成立（引理 7），自同态 \(v_{\varepsilon} = \pi f_{\varepsilon}(u)\) 收敛于 \(\pi f(u)\) ，在 \(\mathcal{L}(\mathrm{E})\) 中（IV, p. 275 的命题 5）。定理得证。

\subsection{预解拓扑与函数演算的连续性}

在本节中，E 表示一个复希尔伯特空间。回顾我们用 \(\mathcal{A}(\mathrm{E})\) 表示 E 上自伴部分算子所成的集。我们将把 IV, p. 194 第 8 目的连续性性质推广到 \(\mathcal{A}(\mathrm{E})\) 上。

定义 6. --- 设 T 是 \(\mathcal{L}(\mathrm{E})\) 上的一个局部凸拓扑。\(\mathcal{A}(\mathrm{E})\) 上的 T-预解拓扑是使映射 \(u \mapsto \mathrm{R}(u, \lambda)\) （从 \(\mathcal{A}(\mathrm{E})\) 到赋以拓扑 T 的 \(\mathcal{L}(\mathrm{E})\) 中）对每个 \(\lambda \in \mathbf{C} - \mathbf{R}\) 都连续的最粗拓扑。

命题 11. --- 设 T 是 \(\mathcal{L}(\mathrm{E})\) 上的局部凸拓扑，它比 \(\mathcal{L}(\mathrm{E})\) 的巴拿赫空间拓扑粗。

\begin{enumerate}
\def\labelenumi{\alph{enumi})}
\item
  设 \(f \in \mathcal{C}_0(\mathbf{R})\) 。对每个序列 \((u_n)_{n \in \mathbf{N}}\) （属于 \(\mathcal{A}(\mathrm{E})\) ，按 \(\mathcal{T}\)-预解拓扑收敛于 u），序列 \((f(u_n))_{n \in \mathbf{N}}\) 收敛于 \(f(u)\) ，在空间 \(\mathcal{L}(\mathrm{E})\) 中，该空间赋以拓扑 \(\mathcal{T}\) ；
\item
  假设每个 \(u \in \mathcal{A}(\mathrm{E})\) 对 \(\mathcal{T}\)-预解拓扑都具有可数的基本邻域系。从 \(\mathcal{C}_0(\mathbf{R}) \times \mathcal{A}(\mathrm{E})\) 到空间 \(\mathcal{L}(\mathrm{E})\) （赋以拓扑 \(\mathcal{T}\) ）中的、由 \((f, u) \mapsto f(u)\) 所定义的映射是连续的。
\end{enumerate}

我们来证明 \(a\) )。拓扑 \(\mathcal{T}\) 在 \(\mathcal{L}(\mathrm{E})\) 上是由对拓扑 \(\mathcal{T}\) 连续的半范数所定义的诸拓扑的上确界（EVT, II, p. 26, 推论及其后的注）。因此，只须证明：对每个半范数 \(p\) ，只要它在 \(\mathcal{L}(\mathrm{E})\) 上对拓扑 \(\mathcal{T}\) 连续，序列 \(p(f(u_n) - f(u))\) 就收敛于 0（TG, I, p. 51, 命题 10）。

设 p 是这样一个半范数。用 \(\mathcal{L}(\mathrm{E})_{p}\) 表示赋以半范数 p 的空间 \(\mathcal{L}(\mathrm{E})\) 的分离巴拿赫空间完备化，并用 \(\varpi\) 表示从 \(\mathcal{L}(\mathrm{E})\) 到 \(\mathcal{L}(\mathrm{E})_{p}\) 中的典范映射；它是连续的，且有 \(p(u) = \|\varpi(u)\|_{p}\) ，对所有 \(u \in \mathcal{L}(\mathrm{E})\) ，其中范数是空间 \(\mathcal{L}(\mathrm{E})_{p}\) 的范数。

由于 T 比 \(\mathcal{L}(\mathrm{E})\) 的巴拿赫空间拓扑粗，存在 \(c \geqslant 0\) 使得 \(p(u) \leqslant c \|u\|\) 对所有 \(u \in \mathcal{L}(\mathrm{E})\) 成立。

先假设 \(f \in \mathcal{D}(\mathbf{R})\) 。设 \(\tilde{f}\) 是 \(f\) 的一个几乎解析延拓。用 \(\mu\) 表示 \(\mathbf{C}\) 上的勒贝格测度。由埃尔弗--舍斯特兰德公式（IV, p. 284 的定理 2），有

\[
\varpi (f (u _ {n})) = - \frac {1}{\pi} \int_ {\mathbf {C}} \frac {\partial \widetilde {f}}{\partial \bar {z}} (\lambda) \varpi (\mathrm{R} (u _ {n}, \lambda)) d \mu (\lambda)\tag{14}
\]

对所有 \(n \in \mathbf{N}\) 成立。

设 \(\lambda \in \mathbf{C} - \mathbf{R}\) 。依据假设，预解式序列 \(\mathrm{R}(u_n,\lambda)\) 收敛于 \(\mathrm{R}(u,\lambda)\) ，在空间 \(\mathcal{L}(\mathrm{E})\) 中，该空间赋以拓扑 \(\mathcal{T}\) ，因此序列 \((\varpi (\mathrm{R}(u_n,\lambda)))_n\) 收敛于 \(\varpi (\mathrm{R}(u,\lambda))\) ，在巴拿赫空间 \(\mathcal{L}(\mathrm{E})_p\) 中。

对每个 \(n \in \mathbf{N}\) ，我们有

\[
\left\| \frac {\partial \widetilde {f}}{\partial \bar {z}} (\lambda) \mathrm{R} (u _ {n}, \lambda) \right\| \leqslant \left| \frac {1}{\mathcal{I} (\lambda)} \frac {\partial \widetilde {f}}{\partial \bar {z}} (\lambda) \right|
\]

（IV, p. 248 的命题 17），因此

\[
\left\| \frac {\partial \widetilde {f}}{\partial \overline {{z}}} (\lambda) \varpi (\mathrm{R} (u _ {n}, \lambda)) \right\| _ {p} \leqslant c \left| \frac {1}{\mathcal {I} (\lambda)} \frac {\partial \widetilde {f}}{\partial \overline {{z}}} (\lambda) \right|.
\]

IV, p. 281 的估计式 (13) 表明，这个不等式的右端对 \(\lambda \in \mathbf{C} - \mathbf{R}\) 是有界函数；它在 \(\mathbf{C}\) 上可积，因为 \(\widetilde{f}\) 具有紧支撑。由勒贝格定理（INT, IV, p. 137, § 3, 第 7 目, 定理 6）以及应用于 \(f\) 的埃尔弗--舍斯特兰德公式，我们推出 \(\varpi(f(u_n))\) 收敛于 \(\varpi(f(u))\) ，从而 \(p(f(u_n) - f(u))\) 趋于 0。

现在假设 \(f \in \mathcal{C}_0(\mathbf{R})\) 。设 \(\varepsilon > 0\) 。存在函数 \(f_{\varepsilon}\) ，它属于 \(\mathcal{D}(\mathbf{R})\) ，使得 \(\| f_{\varepsilon} - f \|_{\infty} \leqslant \varepsilon\) （参见 IV, p. 202 的命题 4, a) 及 INT, III, p. 45, § 1, 第 2 目, 命题 3）。于是，依据 IV, p. 275 的命题 5, c)，有 \(\| f(u) - f_{\varepsilon}(u) \| \leqslant \varepsilon\) 及 \(\| f(u_n) - f_{\varepsilon}(u_n) \| \leqslant \varepsilon\) ，对所有 \(n \in \mathbf{N}\) 。因此，可得

\[
p (f (u _ {n}) - f (u)) \leqslant 2 c \varepsilon + p (f _ {\varepsilon} (u _ {n}) - f _ {\varepsilon} (u))
\]

对所有 \(n \in \mathbf{N}\) 成立。由于 \(f_{\varepsilon} \in \mathcal{D}(\mathbf{R})\) ，由前一种情形，序列 \((p(f_{\varepsilon}(u_n) - f_{\varepsilon}(u)))_{n \in \mathbf{N}}\) 收敛于 0，因此 \(p(f(u_n) - f(u))\) 收敛于 0。这就完成了 a) 的证明。

我们来证明断言 \(b\) )。在关于 \(\mathcal{T}\) 的假设下，由 TG, IX, p. 17 的命题 10 及其后的注，以及断言 \(a\) )，可得映射 \(u \mapsto f(u)\) 从 \(\mathcal{A}(\mathrm{E})\) 到 \(\mathcal{L}(\mathrm{E})\) 中是连续的，当 \(f \in \mathcal{C}_0(\mathbf{R})\) 时。由于映射 \(f \mapsto f(u)\) 从 \(\mathcal{C}_0(\mathbf{R})\) 到 \(\mathcal{L}(\mathrm{E})\) 中连续且范数 \(\leqslant 1\) ，对所有 \(u \in \mathcal{A}(\mathrm{E})\) （IV, p. 275 的命题 5, \(c\) )），断言 \(b\) ) 由 TG, X, p. 13, 推论 3 得出。

例. --- 1) 设 \(\mathscr{S}\) 是 E 的有界子集所成的一个集。\(\mathscr{S}\)-拓扑在 \(\mathcal{L}(\mathrm{E})\) 上（EVT, III, p. 13）是比 \(\mathcal{L}(\mathrm{E})\) 的巴拿赫空间拓扑更粗的局部凸拓扑。

\begin{enumerate}
\def\labelenumi{\arabic{enumi})}
\setcounter{enumi}{1}
\tightlist
\item
  设 \(T_{b}\) 是 \(\mathcal{L}(\mathrm{E})\) 的巴拿赫空间拓扑。对 \(T_{b}\)-预解拓扑，每个 \(u \in \mathcal{A}(\mathrm{E})\) 都具有可数的基本邻域系。
\end{enumerate}

事实上，只须证明：对每个 \(\lambda \in \mathbf{C} - \mathbf{R}\) 及每个 \(\varepsilon > 0\) ，存在 \(\lambda' \in \mathbf{Q} + i\mathbf{Q}^*\) 及整数 \(n \geqslant 1\) ，使得每个部分算子 \(v \in \mathcal{A}(\mathrm{E})\) 若满足 \(\| \mathrm{R}(v, \lambda') - \mathrm{R}(u, \lambda')\| < 1/n\) ，则也满足条件 \(\|R(v,\lambda)-R(u,\lambda)\|<\varepsilon\) ；写出

\[
\begin{array}{c} \| \mathrm{R} (v, \lambda) - \mathrm{R} (u, \lambda) \| \leqslant \| \mathrm{R} (v, \lambda) - \mathrm{R} (v, \lambda^ {\prime}) \| \\ + \| \mathrm{R} (v, \lambda^ {\prime}) - \mathrm{R} (u, \lambda^ {\prime}) \| + \| \mathrm{R} (u, \lambda^ {\prime}) - \mathrm{R} (u, \lambda) \|, \end{array}
\]

这个性质由 IV, p. 245 的公式 (8)、IV, p. 248 的命题 17 的诸估计式，以及 \(\mathbf{Q} + i\mathbf{Q}^*\) 在 \(\mathbf{C} - \mathbf{R}\) 中处处稠密这一事实得出。

若把 \(\mathcal{C}_{0}(\mathbf{R})\) 换成 \(\mathcal{C}_{b}(\mathbf{R})\) ，命题的结论一般不成立（IV, p. 366 习题 29, e)）。尽管如此，我们有下列结果。

推论. --- 设 \(\mathcal{T}_{s}\) 是 \(\mathcal{L}(\mathrm{E})\) 上的逐点收敛拓扑。设 \(f \in \mathcal{C}_{b}(\mathbf{R})\) 。对每个序列 \((u_{n})_{n \in \mathbf{N}}\) （属于 \(\mathcal{A}(\mathrm{E})\) ，收敛于 \(u\) ，按 \(\mathcal{T}_{s}\)-预解拓扑），序列 \((f(u_{n}))_{n \in \mathbf{N}}\) 收敛于 \(f(u)\) ，在 \(\mathcal{L}(\mathrm{E})\) 中，该空间赋以拓扑 \(\mathcal{T}_{s}\) 。

设 \((u_{n})\) 是 \(\mathcal{A}(\mathrm{E})\) 中按 \(T_{s}\)-预解拓扑收敛于 u 的序列。设 \(x \in E\) ，并设 \(\varepsilon > 0\) 。

对每个整数 \(N \in \mathbf{N}\) ，设 \(\varphi_{N} \in \mathcal{K}(\mathbf{R})\) 是支撑包含在 \([-(N+1), N+1]\) 中的函数，使得 \(0 \leqslant \varphi_{N} \leqslant 1\) ，且 \(\varphi_{N}(t) = 1\) 对所有 \(t \in [-N, N]\) 成立。当 N 趋于无穷时，函数 \(\varphi_{N}\) 逐点收敛于 1，且满足 \(|\varphi_{N}| \leqslant 1\) ；因此 IV, p. 276 的命题 6 蕴含存在 \(N \in \mathbf{N}\) 使得 \(\|\varphi_{N}(u)x - x\| \leqslant \varepsilon\) 。置 \(f_{N} = f\varphi_{N}\) 。

对每个 \(n\in \mathbf{N}\) ，我们有

\[
\begin{array}{r l} \| f (u _ {n}) x - f (u) x \| \leqslant & \| f (u _ {n}) x - f _ {\mathrm{N}} (u _ {n}) x \| + \\ & \| f _ {\mathrm{N}} (u _ {n}) x - f _ {\mathrm{N}} (u) x \| + \| f _ {\mathrm{N}} (u) x - f (u) x \|. \end{array} \tag {15}
\]

对每个 \(v\in \mathcal{A}(\mathrm{E})\) ，我们有

\[
\begin{array}{r l} & {\| f (v) x - f _ {\mathrm{N}} (v) x \| = \| (f (1 - \varphi_ {\mathrm{N}})) (v) x \|} \\ & {\qquad \leqslant \| f (v) \|   \| x - \varphi_ {\mathrm{N}} (v) x \| \leqslant \| f \| _ {\infty}   \| x - \varphi_ {\mathrm{N}} (v) x \|} \end{array}
\]

（IV, p. 275 的命题 5, c)）。由于 \(\varphi_{\mathrm{N}} \in \mathcal{C}_0(\mathbf{R})\) ，命题 11 的断言 \(a\) ) 应用于 \(\mathcal{T}_s\)-预解拓扑，蕴含 \(\varphi_{\mathrm{N}}(u_n)x\) 收敛于 \(\varphi_{\mathrm{N}}(u)x\) ，当 \(n \to +\infty\) 时。因此，对所有充分大的 \(n\) ，有

\[
\left\| x - \varphi_ {\mathrm{N}} (u _ {n}) x \right\| \leqslant \left\| x - \varphi_ {\mathrm{N}} (u) x \right\| + \varepsilon \| x \| \leqslant (1 + \| x \|) \varepsilon .
\]

于是不等式 (15) 变成

\[
\| f (u _ {n}) x - f (u) x \| \leqslant \| f \| _ {\infty} (2 + \| x \|) \varepsilon + \| f _ {\mathrm{N}} (u _ {n}) x - f _ {\mathrm{N}} (u) x \|
\]

对所有充分大的 n 成立。对函数 \(f_{\mathrm{N}} \in \mathcal{C}_{0}(\mathbf{R})\) 及 \(T_{s}\)-预解拓扑应用 loc. cit. 的断言 a)，即完成证明。

\subsection{极分解}

引理 8. --- 设 E 是分离的拓扑向量空间。设 \((\mathrm{E}_{1},\|\cdot\|_{1})\) 与 \((\mathrm{E}_{2},\|\cdot\|_{2})\) 是 E 的稠密子空间，其上赋以希尔伯特结构，使得 \(E_{1}\) 与 \(E_{2}\) 到 E 中的典范单射都是连续的。设 F 是 \(E_{1}\cap E_{2}\) 的子空间，它对于这些希尔伯特结构在 \(E_{1}\) 与 \(E_{2}\) 中稠密。若 \(\|x\|_{1}=\|x\|_{2}\) 对所有 \(x\in F\) 成立，则 \(E_{1}=E_{2}\) 且 \(\|\cdot\|_{1}=\|\cdot\|_{2}\) 。

设 \(x \in E_{1}\) 。存在 F 中的序列 \((x_{n})_{n \in \mathbf{N}}\) 使得 \(x_{n}\) 在希尔伯特空间 \(E_{1}\) 中收敛于 x。由于对所有整数 n 与 m 有 \(\|x_{n} - x_{m}\|_{2} = \|x_{n} - x_{m}\|_{1}\) ，序列 \((x_{n})\) 是 \(E_{2}\) 中的柯西序列。用 y 表示它的极限。我们有 \(\|y\|_{2} = \lim\|x_{n}\|_{2} = \|x\|_{1}\) 。由于 \(E_{1}\) 与 \(E_{2}\) 到 E 中的典范单射都是连续的，我们有 \(x_{n} \to x\) 在 E 中成立，同样 \(x_{n} \to y\) 也在 E 中成立。于是由此得出 x = y 且 \(\|x\|_{1} = \|x\|_{2}\) 。特别地，\(E_{1} \subset E_{2}\) 。由对称性即得结论。

设 E 是一个复希尔伯特空间。设 u 是 E 上的正自伴部分算子。对每个 \(\alpha \in \mathbf R_{+}\) ，我们置 \(u^{\alpha} = f(u)\) ，其中 f 是连续映射 \(x \mapsto x^{\alpha}\) ，从 \(\mathbf R_{+}\) 到 \(\mathbf R\) 中。这是一个正自伴部分算子（IV, p. 273 的推论 1, a)）。当 \(u \in \mathcal{L}(E)\) 时，这个记号与关于 C*-代数 \(\mathcal{L}(E)\) 的记号相容（参见 I, p. 118 的命题 16）。若 \(\alpha\) 是正整数，则部分算子 \(u^{\alpha}\) 与由复合 \(u \circ \cdots \circ u\) 所定义的部分算子重合（IV, p. 273 的推论 1, b)）。

设 \(\beta \in \mathbf{R}_{+}\) 。我们有 \(u^{\alpha \beta} = (u^{\alpha})^{\beta}\) （IV, p. 274 的推论 4）。特别地，若 \(\alpha > 0\) ，则部分算子 \(u^{1 / \alpha}\) 是 E 上唯一的正自伴部分算子 \(v\) ，使得 \(v^{\alpha} = u\) 。

假设 \(0 \leqslant \alpha \leqslant \beta\) 。这时我们有 \(\operatorname{dom}(u^{\beta}) \subset \operatorname{dom}(u^{\alpha})\) ：事实上，对每个实数 \(x \geqslant 0\) ，有 \(x^{\alpha} \leqslant 1 + x^{\beta}\) ，于是结论由 \(u^{\alpha}\) 的定义域的定义得出（参见 IV, p. 271 的命题 3）。

设 u 是从 E 到复希尔伯特空间 F 中的闭稠定算子。部分算子 \(u^{*}u\) 是自伴且正的（IV, p. 241 的命题 12）。我们记 \(|u| = (u^{*}u)^{1/2}\) 。

命题 12. --- 设 u 是从 E 到复希尔伯特空间 F 中的闭稠定算子。

\begin{enumerate}
\def\labelenumi{\alph{enumi})}
\item
  正自伴部分算子 \(|u|\) 的定义域与 \(u\) 的定义域重合；
\item
  存在唯一的从 E 到 F 中的部分等距线性映射 \(j\) ，使得 \(u = j|u|\) 且 \(\operatorname{Ker}(j) = \operatorname{Ker}(u)\) ；
\item
  设 \(u_{1}\) 是 E 上的正自伴算子，\(j_{1}\) 是从 E 到 F 中的部分等距线性映射，使得 \(u = j_{1}u_{1}\) 且 \(\operatorname{Ker}(j_{1}) = \operatorname{Ker}(u_{1})\) 。这时我们有 \(u_{1} = |u|\) 且 \(j_{1} = j\) 。
\end{enumerate}

\(u^{*}u\) 的定义域包含在 \(\operatorname{dom}(u)\) 与 \(\operatorname{dom}(|u|)\) 中。它在希尔伯特空间 \(E_{u}\) （loc. cit.）与 \(E_{|u|}\) （IV, p. 273 的推论 1, e)）中稠密，且对每个 \(x \in \operatorname{dom}(u^{*}u)\) ，我们有

\[
\begin{array}{c} \| x \| _ {u} ^ {2} = \| x \| ^ {2} + \langle u (x) | u (x) \rangle = \| x \| ^ {2} + \langle x | (u ^ {*} u) (x) \rangle \\ = \| x \| ^ {2} + \langle | u | (x) | | u | (x) \rangle = \| x \| _ {| u |} ^ {2}. \end{array}\tag{16}
\]

因此，希尔伯特空间 \(E_{u}\) 与 \(E_{|u|}\) 相等（引理 8），从而 \(\text{dom}(u) = \text{dom}(|u|)\) ，且 \(\|u(x)\| = \||u|(x)\|\) 对所有 \(x \in \text{dom}(u)\) 成立。

公式 (16) 蕴含 \(\operatorname{Ker}(u) = \operatorname{Ker}(|u|)\) ，并且存在唯一的等距线性映射 \(v\) ，它从 \(\overline{\operatorname{Im}(|u|)}\) 到 \(\overline{\operatorname{Im}(u)}\) 上，且满足 \(v(|u|(x)) = u(x)\) 对所有 \(x \in \operatorname{dom}(|u|)\) 成立。由于 \(|u|\) 是自伴的，我们有 \(\operatorname{Im}(|u|)^{\circ} = \operatorname{Ker}(|u|)\) （IV, p. 236 的命题 7, c)）。因此存在唯一的从 E 到 F 中的部分等距 \(j\) ，它延拓 \(v\) 且在 \(\operatorname{Ker}(|u|) = \operatorname{Ker}(u)\) 上为零。由于 \(\operatorname{E} = \operatorname{Ker}(u) \oplus \overline{\operatorname{Im}(|u|)}\) ，这个映射是使 \(u = j|u|\) 且 \(\operatorname{Ker}(j) = \operatorname{Ker}(u)\) 的唯一的部分等距映射。

我们来证明 c)。我们有 \(u_{1}j_{1}^{*}j_{1}u_{1}\subset(j_{1}u_{1})^{*}j_{1}u_{1}=u^{*}u\) 。线性映射 \(j_{1}^{*}j_{1}\) 是以 \(\mathrm{Ker}(j_{1})=\mathrm{Ker}(u_{1})\) 为核（EVT, V, p. 41, 命题 5 (ii)），从而以 \(\mathrm{Ker}(u_{1})^{\circ}=\overline{\mathrm{Im}(u_{1})}\) 为像（IV, p. 236 的命题 7, c)）的正交投影算子。因此，\(u_{1}^{2}\subset u^{*}u\) ，又由于这两个算子都是自伴的，故 \(u_{1}^{2}=u^{*}u\) 。于是由此得出 \(u_{1}=(u^{*}u)^{1/2}\) ，而 b) 中的唯一性断言最终表明 \(j_{1}=j\) 。

定义 7. --- 设 u 是从 E 到复希尔伯特空间 F 中的稠定闭算子。由命题 12 所确定的偶 \((j, |u|)\) 称为 u 的极分解。

假设 \(u \in \mathcal{L}(\mathrm{E};\mathrm{F})\) 。它在这个定义意义下的极分解与 I, p. 140 的定义 4 中的极分解重合。

\subsection{由正部分埃尔米特形式定义的自伴算子}

在本节中，E 表示一个复希尔伯特空间。

定义 8. --- 设 D 是 E 的向量子空间。E 上以 D 为定义域的部分埃尔米特形式是指一个对应 \(q = (\Gamma, \mathrm{E} \times \mathrm{E}, \mathbf{C})\) ，其定义域为 \(\mathrm{D} \times \mathrm{D}\) ，其图像 \(\Gamma\) 是函数性的，且从 \(\mathrm{D} \times \mathrm{D}\) 到 \(\mathbf{C}\) 中的、由 \(\Gamma\) 所定义的映射是埃尔米特形式。我们记 \(\operatorname{dom}(q) = \mathrm{D}\) 。

若 \(q\) 所定义的埃尔米特形式是正的，则称它为正部分形式。

设 u 是 E 上的自伴部分算子。对 E 的所有元素 x 与 y，我们用 \(\mu_{x,y}\) （相应地，\(\mu_{x}\) ）表示 \((x,y)\) （相应地，x）关于 u 的谱测度。

设 \((j, |u|)\) 是 \(u\) 的极分解（IV, p. 290 的定义 7）。用 D' 表示 \(|u|^{1/2}\) 的定义域。按定义，它是这样的 \(x \in E\) 所成的空间：函数 \(z \mapsto |z|\) 在 \(\mathrm{Sp}(u)\) 上关于测度 \(\mu_x\) 可积。它包含 \(\mathrm{dom}(u)\) 。

设 \((x,y)\in\mathrm{D}^{\prime}\times\mathrm{D}^{\prime}\) 。\(\mathrm{Sp}(u)\) 的恒等映射关于测度 \(\mu_{x,y}\) 可积（IV, p. 271 的命题 4, b)）；我们置

\[
q _ {u} (x, y) = \int_ {\operatorname{Sp} (u)} t d \mu_ {x, y} (t).
\]

若 \(y \in \operatorname{dom}(u)\) ，则由 IV, p. 271 的公式 (6)，我们有 \(q_u(x, y) = \langle x \mid u(y) \rangle\) 。映射 \(q_u\) 是 D' 上的埃尔米特形式。它定义了 E 上的一个部分埃尔米特形式，称为与 u 相伴的部分埃尔米特形式。若算子 u 是正的，则形式 \(q_u\) 是正部分埃尔米特形式。

定义 9. --- 设 \(q\) 是 E 上的正部分埃尔米特形式。我们用 \(\mathrm{E}_q\) 表示赋以下述纯量积的分离预希尔伯特空间 \(\operatorname{dom}(q)\)

\[
(x \mid y) _ {q} = \langle x \mid y \rangle + q (x, y).
\]

我们用 \(\|x\|_{q}\) 表示 \(x \in E_{q}\) 的范数。若 \(E_{q}\) 是希尔伯特空间，则称形式 q 是闭的。

命题 13. --- 设 u 是 E 上的正自伴部分算子。设 \(q_{u}\) 是与 u 相伴的正部分形式。

\begin{enumerate}
\def\labelenumi{\alph{enumi})}
\item
  \(q_{u}\) 的定义域是 \(u^{1/2}\) 的定义域，且对所有 \(\operatorname{dom}(u^{1/2})\) 中的 x 与 y，我们有 \(q_{u}(x,y)=\langle u^{1/2}(x)\mid u^{1/2}(y)\rangle\) ；
\item
  正部分形式 \(q_{u}\) 是闭的。u 的定义域是 \(q_{u}\) 的核。
\end{enumerate}

由于 \(u\) 是正的，我们有 \(|u| = u\) 。因此 \(\mathrm{dom}(|u|^{1/2})\) （即 \(q_u\) 的定义域）与 \(\mathrm{dom}(u^{1/2})\) 重合，并且我们有

\[
\begin{array}{c} q _ {u} (x, y) = \int_ {\operatorname{Sp} (u)} t d \mu_ {x, y} (t) = \int_ {\operatorname{Sp} (u)} t ^ {1 / 2} \cdot t ^ {1 / 2} d \mu_ {x, y} (t) \\ = \langle u ^ {1 / 2} (x) | u ^ {1 / 2} (y) \rangle \end{array}
\]

对 \(x\) 与 \(y\) （属于 \(\mathrm{dom}(q_u)\) ）成立（IV, p. 271 的命题 4, b)）。这就证明了断言 \(a\) )。

此外，预希尔伯特空间 \(E_{q_{u}}\) 这时与预希尔伯特空间 \(E_{u^{1/2}}\) 重合，后者相伴于 \(u^{1/2}\) （IV, p. 230 的定义 4）。由于 \(u^{1/2}\) 是闭部分算子，这个空间是希尔伯特空间（IV, p. 230 的命题 4），且 \(\text{dom}(u)\) 在 \(E_{u^{1/2}}\) 中稠密（依据 IV, p. 273 的推论 1 的断言 \(e\) )）。

设 \(q\) 是 \(E\) 上的部分埃尔米特形式，其定义域 \(D\) 在 \(E\) 中稠密。对 \(y \in D\) ，用 \(\lambda_y\) 表示线性形式 \(x \mapsto q(y, x)\) （定义在 \(D\) 上）。我们用 \(\widetilde{D}\) 表示这样的 \(y \in D\) 所成的集：线性形式 \(\lambda_y\) 在 \(D\) 上连续。

设 \(y \in \widetilde{D}\) 。由于 D 在 E 中稠密，存在 E 上唯一的连续线性形式，它延拓 \(\lambda_{y}\) 。我们仍用 \(\lambda_{y}\) 表示它。存在 E 中唯一的元素 \(u(y)\) ，使得 \(\lambda_{y}(x) = \langle u(y) | x \rangle\) 对所有 \(x \in E\) 成立（EVT, V, p. 15, 定理 3）。映射 \(y \mapsto u(y)\) 从 \(\widetilde{D}\) 到 E 中是线性的；我们称 E 上以 \(\widetilde{D}\) 为定义域的部分算子 u 为表示 q 的部分算子。因此我们有

\[
q (x, y) = \langle x \mid u (y) \rangle
\]

对 \(y \in \operatorname{dom}(u)\) 及 \(x \in D\) 成立。

注记. --- 设 \(q\) 是闭正部分形式，\(q'\) 是 E 上的连续正埃尔米特形式。在 \(\text{dom}(q)\) 上由 \((x, y) \mapsto q(x, y) + q'(x, y)\) 定义的正埃尔米特形式是以与 \(q\) 相同的定义域为定义域的闭正部分埃尔米特形式。我们把它记作 \(q + q'\) 。

依据 EVT, V, p. 16, 推论 2，存在唯一的线性映射 \(u' \in \mathcal{L}(\mathrm{E})\) ，使得 \(q'(x,y) = \langle x \mid u'(y) \rangle\) 对所有 \((x,y) \in \mathrm{E} \times \mathrm{E}\) 成立。自同态 \(u'\) 是正的，而表示闭正部分埃尔米特形式 \(q + q'\) 的部分算子是 \(u + u'\) 。

命题 14. --- 设 q 是 E 上稠定闭正部分形式，u 是表示 q 的部分算子。

\begin{enumerate}
\def\labelenumi{\alph{enumi})}
\item
  u 的定义域在希尔伯特空间 \(E_{q}\) 中稠密；
\item
  部分算子 u 是自伴且正的。
\end{enumerate}

由于 q 是闭的，定义 9 中的预希尔伯特空间 \(E_{q}\) 是希尔伯特空间。

我们来证明部分算子 \(u + 1_{E}\) （以 \(\operatorname{dom}(u)\) 为定义域）是双射的。由于

\[
\langle x \mid (u + 1 _ {\mathrm{E}}) (x) \rangle = q (x, x) + \| x \| ^ {2} \geqslant \| x \| ^ {2}
\]

对所有 \(x \in \operatorname{dom}(u)\) 成立，这个部分算子是单射的。

设 \(y \in E\) 。\(E_{q}\) 上由 \(x \mapsto \langle y | x \rangle\) 所定义的线性形式是连续的，因为 \(\|x\| \leqslant \|x\|_{q}\) 。因此存在 \(y' \in E_{q}\) 使得

\[
\langle y \mid x \rangle = (y ^ {\prime} \mid x) _ {q} = \langle y ^ {\prime} \mid x \rangle + q (y ^ {\prime}, x)
\]

对所有 \(x \in E_{q}\) 成立（EVT, V, p. 15, 定理 3）。按定义，这意味着 \(y'\) 属于部分算子 \(\widetilde{u}\) 的定义域，该算子表示以 dom(q) 为定义域的部分形式 \((x, y) \mapsto (x \mid y)_{q}\) ，且 \(\widetilde{u}(y') = y\) 。由于由上面的注记有 \(\widetilde{u} = u + 1_{E}\) ，我们由此推出部分算子 \(u + 1_{E}\) 也是满射的，从而是双射的。

我们来证明 \(u\) 的定义域在 \(\mathrm{E}_q\) 中稠密。设 \(y \in \mathrm{E}_q\) 与 \(\operatorname{dom}(u)\) 正交（在 \(\mathrm{E}_q\) 中）。存在 \(y' \in \operatorname{dom}(u)\) 使得 \(u(y') + y' = y\) 。这时我们有

\[
0 = (y \mid y ^ {\prime}) _ {q} = \langle y \mid y ^ {\prime} \rangle + q (y, y ^ {\prime}) = \langle y \mid y ^ {\prime} + u (y ^ {\prime}) \rangle = \| y \| ^ {2}
\]

因此 \(y = 0\) 。于是断言 \(a\) ) 得证。

由于 \(\mathrm{dom}(q)\) 在 E 中稠密，且 \(\| x\| \leqslant \| x\| _q\) 对所有 \(x\in \mathrm{dom}(q)\) 成立，断言 \(a\) ) 蕴含 \(\mathrm{dom}(u)\) 在 E 中稠密。

由于形式 \(q\) 是埃尔米特的（相应地，正的），对所有 \(x\) 与 \(y\) （属于 \(\operatorname{dom}(u)\) ），我们有

\[
\langle y \mid u (x) \rangle = q (y, x) = \overline {{q (x , y)}} = \overline {{\langle x \mid u (y) \rangle}} = \langle u (y) \mid x \rangle ,
\]

（相应地，\(\langle x \mid u(x) \rangle = q(x, x) \geqslant 0\) ），从而 u 是对称的（相应地，正的）。最后，由 IV, p. 240 的命题 10 的推论，部分算子 \(u + 1_{E}\) 是自伴的，u 亦然。

定理 3. --- 使 E 上的正自伴算子 u 对应于正部分形式 \(q_u\) 的映射，是 E 上正自伴部分算子所成的集与 E 上稠定闭正部分形式所成的集之间的一个双射。逆双射使正部分形式 q 对应于表示 q 的部分算子。

由命题 13, b) 与命题 14, b)，陈述中所描述的映射是妥善定义的。我们来证明它们互为逆双射。

设 \(u\) 是 E 上的正自伴部分算子，\(q\) 是与 \(u\) 相伴的正部分形式。用 \(v\) 表示表示 \(q\) 的正自伴算子。设 \(y \in \mathrm{dom}(u) \subset \mathrm{dom}(u^{1/2}) = \mathrm{dom}(q)\) 。对每个 \(x \in \mathrm{dom}(q) = \mathrm{dom}(u^{1/2})\) ，我们有

\[
q (y, x) = \langle u ^ {1 / 2} (y) | u ^ {1 / 2} (x) \rangle = \langle u (y) | x \rangle ,
\]

因此 y 属于 v 的定义域且满足 \(v(y) = u(y)\) 。于是部分算子 v 是 u 的一个延拓；由于 u 与 v 都是自伴的，它们相等。

反之，设 q 是稠定闭正部分形式，u 是表示 q 的正自伴部分算子。我们有 \(\operatorname{dom}(u) \subset \operatorname{dom}(u^{1/2})\) 。对 \(\operatorname{dom}(u)\) 中的 x 与 y，我们有

\[
q _ {u} (x, y) = \langle u ^ {1 / 2} (x) | u ^ {1 / 2} (y) \rangle = \langle x | u (y) \rangle = q (x, y)
\]

（依据命题 13, a)）。于是，希尔伯特空间 \(\mathrm{E}_q\) 与 \(\mathrm{E}_{q_u}\) 都包含 \(\operatorname{dom}(u)\) 作为稠密子空间（分别为命题 14, a) 与命题 13, b)），且 \(\|x\|_q = \|x\|_{q_u}\) 对所有 \(x \in \operatorname{dom}(u)\) 成立。由此得出 \(\mathrm{E}_q = \mathrm{E}_{q_u}\) 且 \(q = q_u\) （引理 8）。

推论. --- 设 \(q\) 是 \(E\) 上的闭正部分形式，而 \(u\) 是表示 \(q\) 的正自伴算子。\(q\) 的定义域等于 \((1_{E} + u)^{1/2}\) 的定义域，且我们有

\[
\| x \| _ {q} = \| (1 _ {\mathrm{E}} + u) ^ {1 / 2} x \|
\]

对所有 \(x \in \operatorname{dom}(q)\) 成立。

\(q\) 的定义域等于 \(u^{1/2}\) 的定义域（命题 13, \(a\) )），后者与 \((1_{\mathrm{E}} + u)^{1/2}\) 的定义域重合（参见 IV, p. 271 的命题 3）。对每个 \(x \in \operatorname{dom}(u) \subset \operatorname{dom}(u^{1/2})\) ，我们有

\[
\| (1 _ {\mathrm{E}} + u) ^ {1 / 2} x \| ^ {2} = \langle x | (1 _ {\mathrm{E}} + u) (x) \rangle = \| x \| ^ {2} + \langle x | u (x) \rangle = \| x \| _ {q} ^ {2}.
\]

由于 \(u\) 的定义域在希尔伯特空间 \(\mathrm{E}_q\) 中稠密（命题 14, a)），这个公式由连续性延拓到所有 \(x \in \mathrm{dom}(u^{1/2})\) 。

例. --- 设 u 是 E 上的正部分算子，它未必是闭的。用下式定义以 dom(u) 为定义域的正部分形式 q

\[
q (x, y) = \langle x \mid u (y) \rangle
\]

对 \(x\) 与 \(y\) （属于 \(\mathrm{dom}(u)\) ）成立。设 \(\mathrm{E}_q\) 是 IV, p. 292 的定义 9 中的分离预希尔伯特空间。设 \(\widetilde{\mathrm{E}}_q\) 是 \(\mathrm{E}_q\) 的完备化希尔伯特空间，其纯量积仍记作 \((x,y) \mapsto (x \mid y)_q\) 。由于 \(\iota\) （\(\mathrm{E}_q\) 到 E 中的典范包含）是连续的，它具有唯一的连续延拓，记作 \(\widetilde{\iota}\) （从 \(\widetilde{\mathrm{E}}_q\) 到 E 中）。埃尔米特形式 \(q\) 在 \(\mathrm{E}_q\) 上是连续的，因而也延拓为唯一的连续正埃尔米特形式 \(\widetilde{q}\) ，定义在 \(\widetilde{\mathrm{E}}_q\) 上。

我们来证明线性映射 \(\widetilde{\iota}\) 是单射的。设 \(x\in\mathrm{Ker}(\widetilde{\iota})\) 。考虑序列 \((x_{n})_{n\in\mathbf{N}}\) （属于 \(\mathrm{E}_{q}\) ），它收敛于 \(x\) ，在 \(\widetilde{\mathrm{E}}_{q}\) 中。这时我们有

\[
\lim _ {n \to + \infty} \iota (x _ {n}) = \widetilde {\iota} (x) = 0,
\]

因此序列 \((x_{n})_{n\in\mathbf{N}}\) 在 E 中收敛于 0。设 \(y\in E_{q}\) 。由于 \(\widetilde{q}\) 在 \(\widetilde{E}_{q}\) 上连续，由此可得

\[
\begin{array}{r l} (x \mid y) _ {q} = \lim _ {n \to + \infty} (x _ {n} \mid y) _ {q} = & \lim _ {n \to + \infty} \Big (\langle x _ {n} \mid y \rangle + q (x _ {n}, y) \Big) \\ & = \lim _ {n \to + \infty} \Big (\langle x _ {n} \mid y \rangle + \langle x _ {n} \mid u (y) \rangle \Big) = 0. \end{array}
\]

由于空间 \(E_{q}\) 在 \(\widetilde{E}_{q}\) 中稠密，我们推出 x = 0，此即所要证者。

把 \(\widetilde{E}_{q}\) 与它在 \(\widetilde{\iota}\) 下的像等同于 E 的子空间，我们便把 \(\widetilde{q}\) 解释为延拓 q 的稠定闭正部分形式。与 \(\widetilde{q}\) 相伴的正自伴算子是 u 的一个自伴延拓。它称为 u 的弗里德里希斯扩张。

注. --- 设 \(c \in \mathbf R_{+}\) 。设 q 是稠定闭部分埃尔米特形式，使得对所有属于 q 的定义域的 x 有 \(q(x, x) \geqslant -c \|x\|^2\) 。这意味着由 \(\widetilde{q}(x, y) = q(x, y) + c \langle x \mid y \rangle\) 定义的以 dom(q) 为定义域的闭部分埃尔米特形式是正部分形式。于是由定理 3，这样的形式对应于 E 上使 \(u + c1_{E}\) 为正的自伴部分算子 u，亦即使得 u 以 -c 为下界（IV, p. 237 的定义 7）者。

反之，设 u 是 E 上的对称部分算子，使得

\[
\langle x \mid u (x) \rangle \geqslant - c \| x \| ^ {2}\tag{17}
\]

对所有 \(x \in \operatorname{dom}(u)\) 成立。这时 \(v = u + c1_{E}\) 是 E 上的正部分算子。u 的弗里德里希斯扩张是自伴算子 \(\widetilde{v} - c1_{E}\) ，其中 \(\widetilde{v}\) 是 v 的弗里德里希斯扩张；它是 u 的一个自伴延拓，不依赖于满足 (17) 的实数 c 的选取。